%% 熱-8 水蒸気で回す：蒸気表とランキンサイクル（第22回編集委員会 物理第2巻：組 熱W5）
%% 記号と言い方（計画書 C-1・B の熱-8）：
%%   圧力 p（絶対圧。ゲージ圧はことばで書きそえる）、絶対温度 T〔K〕・セルシウス温度 θ〔℃〕（第22回 2章の4。表の温度の列は θ〔℃〕。整合の組で直した）、
%%   比体積 v〔m^3/kg〕（この章では速さを使わない）、比エンタルピー h、比エントロピー s（熱-6）、乾き度 x、
%%   飽和液の量に ′、飽和蒸気の量に ″（h′・h″・v′・v″・s′・s″。JIS・日本の蒸気表の書き方。英語の本では h_f・h_g）、
%%   蒸発熱 L＝h″−h′（熱-1 の L と同じもの）。
%%   サイクルの熱と仕事は、大きさ（正の数）と向きの添字で書く（第22回 2章の3）：ボイラーで受けとる熱 q_入、
%%   復水器で捨てる熱 q_出、タービンが取り出す仕事 w_出、ポンプに入れる仕事 w_入、正味の仕事は w_出−w_入（T-1 で w_正味 は使わない）。
%%   効率 η、質量流量 m-dot、電気の出力 P。
%% 数値：checks/物理第2巻/熱-8/calc_p2_08.py（出力 calc_p2_08_out.txt、状態点の一覧 states_p2_08.csv）。
%%   水・水蒸気は IAPWS-IF97 の式（iapws 1.5.5）。CoolProp 8.0（IAPWS-95）との差は飽和の h で最大 0.2 kJ/kg。
%%   本文の表の丸め：温度は小数2けた、h は小数1けた、s は小数4けた、v′ は小数6けた、v″ は4けた。巻末の蒸気表も同じ丸めにする（報告）。
%% 出典・確かめたこと：
%%   三重点 273.16 K・611.657 Pa、臨界点 647.096 K・22.064 MPa：IAPWS-IF97（改訂版）の定数（iapws の値で確認）。
%%   1気圧の沸点 99.97 ℃：IF97 の飽和線（ITS-90 の温度）。昔の温度目盛（IPTS-68 まで）では 100 ℃と決めていた。
%%   モリエ線図：R. Mollier, "Neue Diagramme zur technischen Wärmelehre", Z. VDI 48 (1904) 271–275。
%%   ランキン（W. J. M. Rankine）：『A Manual of the Steam Engine and Other Prime Movers』1859年。
%%   ワットの分離凝縮器の特許 1769年。パーソンズの反動タービン 1884年。
%%   A-USC（先進超々臨界圧）：経済産業省の技術評価資料（平成25年、電力分野）・NEDO の資料。蒸気温度700 ℃級で送電端効率46 %（高位発熱量基準）を目標に開発。
%%   最新の超々臨界圧（USC）の送電端効率 約42 %（高位発熱量基準）：同じ資料。先進12Cr鋼の限界が約630 ℃、その上はニッケル基合金。
%%   沸騰水型原子炉の主蒸気 約7 MPa・約286 ℃（飽和蒸気）、復水器の圧力 約5 kPa：発電の教科書の典型値（〔要確認〕の値は報告に書いた）。
%%   エベレスト山頂の気圧 約34 kPa（West ほか 1983年の測定 253 Torr＝33.7 kPa）。
\makeatletter
% 飽和線のドーム（T-s）と、相図などで使う線の色
\tikzset{qxi lab/.style={font=\sffamily\footnotesize, inner sep=1pt},
  qxi arr/.style={-{Stealth[length=3.6pt, width=3.2pt]}, line width=1pt},
  qxi flow/.style={-{Stealth[length=4.2pt, width=3.8pt]}, line width=1.4pt, bkBlue!80!black},
  qxi heat/.style={-{Stealth[length=4.2pt, width=3.8pt]}, line width=1.4pt, bkRed!80!black},
  qxi work/.style={-{Stealth[length=4.2pt, width=3.8pt]}, line width=1.4pt, bkOrange!85!black},
  qxi box/.style={draw=bkInk, line width=0.8pt, fill=bkGray!12, rounded corners=1.2mm, font=\sffamily\footnotesize, align=center},
  qxi num/.style={circle, draw=bkInk, fill=white, line width=0.6pt, inner sep=0.6pt, font=\sffamily\footnotesize\bfseries}}
% 状態の番号の札  \qxinum{1}
\newcommand\qxinum[1]{\tikz[baseline=(n.base)]\node[qxi num](n){#1};}
% ランキンサイクルの装置の図（ボイラー・タービン・復水器・ポンプ）
\newcommand\qxirankine{%
\begin{tikzpicture}[x=1mm, y=1mm]
  % ボイラー
  \filldraw[fill=bkRed!8, draw=bkInk, line width=0.8pt, rounded corners=1.5mm] (0,0) rectangle (26,34);
  \node[qxi lab, font=\sffamily\small\bfseries] at (13,30) {ボイラー};
  \foreach \y in {9,14,19,24} {\draw[bkInk!70, line width=0.6pt] (5,\y) -- (21,\y);}
  \foreach \x in {6,11,16,21} {\fill[bkOrange!80] (\x-1.6,1.5) .. controls (\x-1.6,4) and (\x,4.5) .. (\x,6) .. controls (\x,4.5) and (\x+1.6,4) .. (\x+1.6,1.5) -- cycle;}
  \draw[qxi heat] (13,-11) -- (13,-0.8);
  \node[qxi lab, anchor=north] at (13,-11.8) {$q_{\text{入}}$（燃料の熱）};
  % タービン・発電機
  \filldraw[fill=bkBlue!10, draw=bkInk, line width=0.8pt] (60,30) -- (60,46) -- (84,52) -- (84,24) -- cycle;
  \node[qxi lab, font=\sffamily\small\bfseries] at (72,38) {タービン};
  \draw[bkInk, line width=1.6pt] (84,38) -- (92,38);
  \filldraw[fill=bkGray!25, draw=bkInk, line width=0.8pt, rounded corners=1mm] (92,31) rectangle (110,45);
  \node[qxi lab] at (101,38) {発電機};
  \draw[qxi work] (110.8,38) -- (122,38);
  \node[qxi lab, anchor=south] at (116.5,39) {$w_{\text{出}}$};
  % 復水器
  \filldraw[fill=bkBlue!6, draw=bkInk, line width=0.8pt, rounded corners=1.5mm] (66,-4) rectangle (98,12);
  \node[qxi lab, font=\sffamily\small\bfseries] at (82,8.3) {復水器};
  \draw[bkBlue!70!black, line width=0.9pt, rounded corners=0.8mm] (102,0) -- (72,0) -- (72,3.5) -- (102,3.5);
  \draw[qxi heat] (103,1.8) -- (116,1.8);
  \node[qxi lab, anchor=north east] at (122,-1.4) {$q_{\text{出}}$（海の水へ）};
  % ポンプ
  \filldraw[fill=bkGray!20, draw=bkInk, line width=0.8pt] (46,4) circle (5.2);
  \draw[bkInk, line width=0.8pt] (43,0.5) -- (46,7.5) -- (49,0.5);
  \node[qxi lab, anchor=south] at (46,9.8) {ポンプ};
  \draw[qxi work] (46,-12) -- (46,-1.6);
  \node[qxi lab, anchor=west] at (47.5,-9) {$w_{\text{入}}$};
  % 流れ
  \draw[qxi flow] (26,30) -- (40,30) -- (40,42) -- (59.2,42);
  \draw[qxi flow] (80,23.2) -- (80,12.8);
  \draw[qxi flow] (66,4) -- (51.8,4);
  \draw[qxi flow] (40.2,4) -- (26.8,4);
  \node[qxi num] at (47,46) {1};
  \node[qxi num] at (84.5,18) {2};
  \node[qxi num] at (58.5,8.2) {3};
  \node[qxi num] at (33,8.2) {4};
\end{tikzpicture}}
\makeatother

\phchapter{8}{水蒸気で回す：\\蒸気表とランキンサイクル}{高い山の上では、なぜ水が100\,℃より低い温度でわく？}{\bklevelrange{2}{4}}{水蒸気で回す：蒸気表とランキンサイクル}
\bkmichishirube{3-5,3-12,熱-1,熱-4,熱-5,熱-6}{熱-9,熱-10,熱-16}

\begin{goal}
\begin{itemize}[label={\color{bkTheme}●}, leftmargin=1.6\zw, itemsep=0pt, topsep=1pt]
\item \bklevel{2}\enspace 水のわく温度が、周りの圧力で決まることを、飽和蒸気圧の表で説明できる。
\item \bklevel{4}\enspace 乾き度と蒸気表（直線補間）で、湿り蒸気・過熱蒸気の比エンタルピーを求められる。
\item \bklevel{4}\enspace ランキンサイクルの効率を蒸気表から計算し、再熱・再生の働きを説明できる。
\end{itemize}
\end{goal}

%% 「この章で導く中学・高校の公式」の欄は置かない（第23回 U-5・K-8。もとの行は小学校理科の「わかること」で、公式ではないのでけずった）。

\begin{hatena}
\noindent\begin{minipage}[c]{0.6\linewidth}\raggedright
富士山の頂上でお湯をわかすと、ぐらぐらわいているのに、温度計は約\qty{87}{\celsius}しか示しません。カップめんは、少しかために仕上がります。

火を強くしても、わいている水の温度は上がりません。わく温度を決めているのは、火の強さではなく何なのでしょう。
\end{minipage}\hfill
\begin{minipage}[c]{0.36\linewidth}\centering
\begin{tikzpicture}[x=1mm, y=1mm]
  \fill[bkGray!25] (-4,0) -- (16,30) -- (22,26) -- (40,0) -- cycle;
  \fill[white] (11.5,23.2) -- (16,30) -- (22,26) -- (25.3,21.3) -- (21,22.6) -- (18,21) -- (15,23) -- cycle;
  \draw[bkInk, line width=0.7pt] (-4,0) -- (16,30) -- (22,26) -- (40,0);
  \filldraw[fill=bkGray!45, draw=bkInk, line width=0.7pt] (10,31) rectangle (24,38);
  \fill[bkBlue!30] (10.6,31.6) rectangle (23.4,36);
  \foreach \x in {13,16,22} {\draw[bkBlue!70!black, line width=0.5pt] (\x,34.6) circle (0.8);}
  \draw[bkInk, line width=0.7pt] (19,33) -- (19,48);
  \fill[bkRed!75!black] (19,33) circle (1.3);
  \draw[bkRed!75!black, line width=1.2pt] (19,33) -- (19,42);
  \node[qxi lab, anchor=west] at (21,45) {約\qty{87}{\celsius}};
\end{tikzpicture}
\end{minipage}
\end{hatena}

\noindent 熱-1 では、水がわいている間、加えた熱は温度を上げずに、水を水蒸気に変えるのに使われることを学びました（蒸発熱）。この章では、まず、わく温度そのものが\textbf{周りの圧力}で決まることを調べます。次に、わいている途中の水や、わいたあとの水蒸気の「熱の量」を、\textbf{蒸気表}という表から読む方法を学びます。最後に、その表を使って、火力発電所や原子力発電所の心臓部である\textbf{ランキンサイクル}の効率を計算します。熱-0 の発電所で、燃料の熱の半分以上を海へ捨てていたわけも、ここで数で確かめます。

%% ---------------- やってみよう・思い出そう ----------------
\begin{yattemiyou}[注射器の中で、ぬるま湯をわかそう]
\textbf{用意するもの}\par
先にふたができるプラスチックの注射器（\qty{20}{mL}くらい。先をふさぐゴムのキャップ）、ぬるま湯（\qtyrange{40}{50}{\celsius}。やけどしない温度）、温度計
\begin{enumerate}[label={\bkcircled[bkGreen]{\arabic*}}, leftmargin=2.4\zw, itemsep=1pt, topsep=2pt]
\item ぬるま湯の温度を測り、注射器に\qty{3}{mL}くらい吸いこむ。中の空気はできるだけ押し出す。
\item 先をキャップでふさぎ、ピストンを強く引いて止める。湯のようすを見る。
\item ピストンをもどすと、どうなるかも見る。湯の温度を変えて、わき方をくらべる。
\end{enumerate}
\end{yattemiyou}

\noindent ピストンを引くと、ぬるま湯から、泡がぶくぶくと出てきます。中の圧力を下げると、\qty{100}{\celsius}よりずっと低い温度でも水がわくのです。熱いお湯を使うと危ないので、必ずぬるま湯で行いましょう。なぜこうなるかは、①で説明します。

\begin{omoidasou}
\bkrecall{熱-1}{わいている間は温度が一定で、加えた熱は蒸発に使われる。1 kg を蒸発させる熱が蒸発熱 $L$（\qty{100}{\celsius}で約\qty{2.26}{MJ/kg}）。}
\bkrecall{熱-4}{流れる物の1 kg あたりのエネルギーのバランス：$q+w=h_2-h_1$（$q$ は受けとる熱、$w$ はされる仕事。速さ・高さの差がないとき）。取り出す仕事を $w_{\text{出}}=-w$ と書けば $q_{\text{入}}-w_{\text{出}}=h_2-h_1$。$h$ は比エンタルピー。}
\bkrecall{熱-5}{熱機関の効率は $\eta=W_{\text{出}}/Q_{\mathrm H}$。上限はカルノーの効率 $1-T_{\mathrm L}/T_{\mathrm H}$（温度は K）。}
\bkrecall{熱-6}{エントロピー $\Delta S=Q/T$。断熱で可逆な変化では $s$ は一定。$T$-$s$ 線図で、線の下の面積が熱。}
\bkrecall{3-5}{表のあいだの値を、2つの点を結ぶ直線で読む（直線内挿）。}
\bkrecall{3-12}{倍々に増える量は、片対数グラフで直線に見える。}
\end{omoidasou}

%% ---------------- 物理の道具 ----------------
\bksection{doguu}{物理の道具}

\bksubsection{1}{わく温度は、圧力で決まる：飽和蒸気圧}

\leavevmode\label{ph:熱-8:sec1}コップの水は、わかさなくても、少しずつ減っていきます（蒸発）。水の表面では、速く動いている分子が、周りの分子を振り切って空気の中へ飛び出しているからです（分子の速さのばらつきは熱-14）。反対に、空気の中の水の分子が、水面にぶつかって水にもどることもあります。

\noindent\begin{minipage}[c]{0.40\linewidth}\centering
\begin{tikzpicture}[x=1mm, y=1mm]
  \draw[bkInk, line width=0.9pt] (0,40) -- (0,0) -- (40,0) -- (40,40);
  \fill[bkBlue!20] (0.5,0.5) rectangle (39.5,16);
  \draw[bkBlue!70!black, line width=0.7pt] (0.5,16) -- (39.5,16);
  \filldraw[fill=bkGray!40, draw=bkInk, line width=0.8pt] (-1,40) rectangle (41,43);
  \foreach \p in {(5,5),(12,9),(20,4),(28,10),(34,5),(8,13),(24,13),(16,11),(32,13.5)} {\fill[bkBlue!70!black] \p circle (1.1);}
  \foreach \p in {(5,35),(20,36),(34,33),(14,30),(26,31)} {\fill[bkBlue!50] \p circle (1.1);}
  \draw[qxi arr, bkRed!75!black] (11,16.6) -- (9,22.5);
  \draw[qxi arr, bkBlue!80!black] (29,23) -- (28,16.8);
  \node[qxi lab, anchor=south] at (9,23) {飛び出す};
  \node[qxi lab, anchor=south] at (29.5,23.5) {もどる};
\end{tikzpicture}
\end{minipage}\hfill
\begin{minipage}[c]{0.56\linewidth}\raggedright
ふたをした容器に水を入れて、温度を一定に保ちます。はじめは飛び出す分子のほうが多いのですが、空気の中の水蒸気が増えると、もどる分子も増えます。やがて、飛び出す数ともどる数が同じになり、水蒸気の量はそれ以上増えなくなります。この状態を\textbf{飽和}といい、そのときの水蒸気の圧力を\textbf{飽和蒸気圧} $p_{\mathrm s}$ といいます。\bkidx{ほうわじょうきあつ}{飽和蒸気圧}
\end{minipage}
\par\vspace{3pt}
温度を上げると、飛び出す分子が増えるので、飽和蒸気圧は大きくなります。飽和蒸気圧は\textbf{温度だけで決まり}、容器の大きさや、ほかの空気が混ざっているかにはよりません。表\ref{ph:熱-8:tab-ps}は、水の飽和蒸気圧です。力-14 で使った「\qty{20}{\celsius}の水の飽和蒸気圧は約\qty{2.3}{kPa}」は、この表の\qty{2.339}{kPa}です。

\begin{bkcenter}
\begin{bktable}{colspec={Q[c,10mm] Q[r,18mm] Q[r,18mm] | Q[c,10mm] Q[r,18mm] Q[r,18mm]}, row{1}={font=\sffamily\footnotesize\bfseries}}
$\theta$〔℃〕 & $p_{\mathrm s}$〔kPa〕 & $L$〔kJ/kg〕 & $\theta$〔℃〕 & $p_{\mathrm s}$〔kPa〕 & $L$〔kJ/kg〕 \\
0.01 & 0.6117 & 2500.9 & 80 & 47.415 & 2308.1 \\
10 & 1.2282 & 2477.2 & 100 & 101.42 & 2256.5 \\
20 & 2.3392 & 2453.5 & 120 & 198.67 & 2202.1 \\
30 & 4.2467 & 2429.8 & 150 & 476.10 & 2113.7 \\
40 & 7.3844 & 2406.0 & 200 & 1555 & 1939.7 \\
50 & 12.351 & 2382.0 & 250 & 3976 & 1715.3 \\
60 & 19.946 & 2357.7 & 300 & 8588 & 1404.8 \\
\end{bktable}
\bktabcaption{水の飽和蒸気圧 $p_{\mathrm s}$ と蒸発熱 $L$（IAPWS-IF97 の式から計算。巻末の蒸気表の抜き出し）。}
\label{ph:熱-8:tab-ps}
\end{bkcenter}

\noindent\textbf{わく（沸騰する）とは}\par\nobreak
水を下から温めると、底のあたりで水の中に水蒸気の泡ができようとします。泡の中の圧力は、その温度の飽和蒸気圧 $p_{\mathrm s}$ です。泡の周りは、水面を押す大気の圧力 $p$ で押されています（深さの分の水の圧力は小さいので、ここでは無視します）。$p_{\mathrm s}$ が $p$ より小さいうちは、泡はつぶされてしまいます。温度が上がって $p_{\mathrm s}$ が $p$ に届くと、泡はつぶれずに育ち、水の中からさかんに水蒸気が出てきます。これが\textbf{沸騰}です。\bkidx{ふっとう}{沸騰}
\pheq{p_{\mathrm s}(T)=p\quad\text{（わく温度 $T$ の決め方）}}{kPa ＝ kPa}
\label{ph:熱-8:futten}1気圧（\qty{101.325}{kPa}）で $p_{\mathrm s}$ がちょうど1気圧になる温度は\qty{99.97}{\celsius}で、これが「水は約\qty{100}{\celsius}でわく」の正体です。わいている間に火を強くしても、$p_{\mathrm s}=p$ の温度より上がることはできません。加えた熱は、泡を増やす（水を水蒸気に変える）のに使われます。だから、わいている間は温度が変わらないのです（熱-1）。\bkidx{ふってん}{沸点}

\begin{phyomi}[$p_{\mathrm s}(T)=p$]
「わく温度は、飽和蒸気圧が周りの圧力と等しくなる温度」。周りの圧力 $p$ が決まれば、表\ref{ph:熱-8:tab-ps}を逆に読んで、わく温度が決まります。圧力を下げれば低い温度で、上げれば高い温度でわきます。
\end{phyomi}

\bkhatenaref の答えです。富士山の頂上の気圧は約\qty{63}{kPa}です（力-12）。表\ref{ph:熱-8:tab-ps}では、$p_{\mathrm s}$ が\qty{63}{kPa}になる温度は、\qty{80}{\celsius}（\qty{47.4}{kPa}）と\qty{100}{\celsius}（\qty{101.4}{kPa}）の間にあり、詳しく計算すると約\qty{87}{\celsius}です（\bkreidairef{1}）。周りの圧力が低いので、水は\qty{100}{\celsius}になる前に、わいてしまうのです。やってみようの注射器も同じで、ピストンを引いて中の圧力を\qty{40}{\celsius}の $p_{\mathrm s}$（約\qty{7.4}{kPa}）より下げたので、ぬるま湯がわきました。

\begin{phmijika}[圧力なべ・井戸のポンプ・キャビテーション]
\textbf{圧力なべ}は、ふたで中の水蒸気を閉じこめて、圧力を約2気圧に上げます。表\ref{ph:熱-8:tab-ps}では、\qty{120}{\celsius}で $p_{\mathrm s}$ が約\qty{199}{kPa}なので、中の水は約\qty{120}{\celsius}でわきます（\bkreidairef{13}）。高い温度で煮るので、早く煮えるのです。
力-12 では、井戸の水を吸い上げるポンプが、どうしても約\qty{10}{m}より高くは水を上げられないことを学びました。吸いこみ口の圧力を下げすぎると、まず水がその温度でわいて（$p_{\mathrm s}$ より下がって）、水蒸気の泡ができてしまうからでもあります。ポンプの羽根車の周りで起きるこの現象が、力-14 のキャビテーションです。
\end{phmijika}

\begin{phbaisuru}[\qty{20}{K}高い温度]{温度 $T$}
飽和蒸気圧は、温度とともに急に増えます。\qty{20}{\celsius}から\qty{40}{\celsius}で約3.2倍、\qty{40}{\celsius}から\qty{60}{\celsius}で約2.7倍、\qty{60}{\celsius}から\qty{80}{\celsius}で約2.4倍、\qty{80}{\celsius}から\qty{100}{\celsius}で約2.1倍です。\qty{20}{K}ごとに「何倍か」になる、指数関数（3-11）に近い増え方です。だから、片対数グラフ（3-12⑥）にかくと、ほぼ直線になります（詳しい式は熱-10）。
\par\vspace{2pt}\noindent
\begin{minipage}[t]{0.49\linewidth}\centering
{\sffamily\footnotesize ふつうの目もり}\par
\begin{tikzpicture}
\begin{axis}[bkplot, width=46mm, height=42mm, xmin=0, xmax=104, ymin=0, ymax=110, xtick={0,50,100}, ytick={0,50,100}, xlabel={温度（℃）}, ylabel={$p_{\mathrm s}$（kPa）}]
  \addplot[bkcurve, smooth] coordinates {(0,0.611) (5,0.873) (10,1.228) (15,1.706) (20,2.339) (25,3.170) (30,4.247) (35,5.629) (40,7.384) (45,9.594) (50,12.351) (55,15.761) (60,19.946) (65,25.041) (70,31.201) (75,38.595) (80,47.415) (85,57.867) (90,70.182) (95,84.609) (100,101.418)};
  \addplot[only marks, mark=*, mark size=1.6pt, bkOrange] coordinates {(20,2.339) (40,7.384) (60,19.946) (80,47.415) (100,101.418)};
\end{axis}
\end{tikzpicture}
\end{minipage}\hfill
\begin{minipage}[t]{0.49\linewidth}\centering
{\sffamily\footnotesize 縦を対数の目もりにしたグラフ}\par
\begin{tikzpicture}
\begin{axis}[bkplot, axis lines=left, width=46mm, height=42mm, xmin=0, xmax=104, ymode=log, ymin=0.5, ymax=150, xtick={0,50,100}, ytick={1,10,100}, yticklabels={1,10,100}, xlabel={温度（℃）}, ylabel={$p_{\mathrm s}$（kPa）}]
  \addplot[bkcurve, smooth] coordinates {(0,0.611) (10,1.228) (20,2.339) (30,4.247) (40,7.384) (50,12.351) (60,19.946) (70,31.201) (80,47.415) (90,70.182) (100,101.418)};
  \addplot[only marks, mark=*, mark size=1.6pt, bkOrange] coordinates {(20,2.339) (40,7.384) (60,19.946) (80,47.415) (100,101.418)};
\end{axis}
\end{tikzpicture}
\end{minipage}
\par\raggedright\footnotesize 対数の目もりでは、ほぼまっすぐになる（少し上に凸）。
\end{phbaisuru}

\noindent\textbf{水の相図：三重点と臨界点}\par\nobreak
温度と圧力の平面に、水が氷・水・水蒸気のどれになるかをかきこんだ図を\textbf{相図}といいます（図\ref{ph:熱-8:fig-sozu}）。\bkidx{そうず}{相図}水と水蒸気の境の線が、表\ref{ph:熱-8:tab-ps}の飽和蒸気圧の線です。この線の上では、水と水蒸気がいっしょにいられます。
\begin{bkfigure}
\begin{tikzpicture}
\begin{axis}[bkplot, axis lines=left, width=108mm, height=72mm, xmin=-40, xmax=420, ymode=log, ymin=0.01, ymax=2e5, xtick={0,100,200,300,400}, ytick={0.01,0.1,1,10,100,1000,10000,100000}, yticklabels={0.01,0.1,1,10,100,1000,\num{10000},\num{100000}}, xlabel={温度（℃）}, ylabel={圧力（kPa）}]
  \addplot[bkcurve, smooth] coordinates {(0.01,0.6117) (20,2.339) (50,12.35) (100,101.4) (150,476.1) (200,1555) (250,3976) (300,8588) (350,16529) (373.95,22064)};
  \addplot[bkcurve2, densely dashed, smooth] coordinates {(-40,0.0129) (-30,0.0381) (-20,0.1032) (-10,0.2599) (0.01,0.6117)};
  \addplot[bkcurve2, densely dashed] coordinates {(0.01,0.6117) (-1.15,15136) (-9,100084)};
  \addplot[only marks, mark=*, mark size=2pt, bkInk] coordinates {(0.01,0.6117) (373.95,22064)};
  \addplot[only marks, mark=o, mark size=2pt, bkOrange!80!black, line width=0.9pt] coordinates {(99.97,101.325) (87.18,63) (120.21,200)};
  \draw[bkGray, line width=0.4pt] (axis cs:121.5,190) -- (axis cs:128,150);
  \draw[bkGray, line width=0.4pt] (axis cs:101,95) -- (axis cs:106,62);
  \draw[bkGray, line width=0.4pt] (axis cs:88.5,58) -- (axis cs:92,24);
  \node[font=\sffamily\footnotesize, anchor=west] at (axis cs:128,150) {圧力なべ};
  \node[font=\sffamily\footnotesize, anchor=west] at (axis cs:106,60) {1気圧};
  \node[font=\sffamily\footnotesize, anchor=west] at (axis cs:92,22) {富士山頂};
  \node[font=\sffamily\footnotesize, anchor=west] at (axis cs:14,0.35) {三重点（\qty{0.01}{\celsius}・\qty{0.61}{kPa}）};
  \node[font=\sffamily\footnotesize, anchor=south east] at (axis cs:368,30000) {臨界点（\qty{374}{\celsius}・\qty{22.1}{MPa}）};
  \node[font=\sffamily\small\bfseries] at (axis cs:-24,1000) {氷};
  \node[font=\sffamily\small\bfseries] at (axis cs:130,8000) {水（液体）};
  \node[font=\sffamily\small\bfseries] at (axis cs:270,5) {水蒸気（気体）};
\end{axis}
\end{tikzpicture}
\bkfigcaption{水の相図（縦は対数の目もり）。実線が水と水蒸気の境（飽和蒸気圧の線）、破線が氷との境。白い丸は、1気圧（\qty{101.3}{kPa}・\qty{99.97}{\celsius}）、富士山頂（\qty{63}{kPa}・約\qty{87}{\celsius}）、圧力なべ（\qty{200}{kPa}・約\qty{120}{\celsius}）。}
\label{ph:熱-8:fig-sozu}
\end{bkfigure}
3本の線が出会う点を\textbf{三重点}といい、氷・水・水蒸気の3つがいっしょにいられる、ただ1つの温度と圧力（\qty{0.01}{\celsius}・\qty{611.657}{Pa}）です。\bkidx{さんじゅうてん}{三重点}水と水蒸気の線は、上の方の\textbf{臨界点}（\qty{373.946}{\celsius}・\qty{22.064}{MPa}）で終わります。\bkidx{りんかいてん}{臨界点}臨界点より高い温度と圧力では、水と水蒸気の区別がなくなり、「わく」ことはありません。最新の火力発電所は、この臨界点より高い圧力で水を温めています（⑧）。

\begin{seikaku}
氷と水の境の線が左へ傾いているのは、水が氷になると体積が増える（小学校理科）からです。氷を強く押すと、体積の小さい水のほうになろうとして、0\,℃より少し低い温度でもとけます。ただし、その下がり方は\qty{10}{MPa}（約100気圧）押しても約\qty{0.7}{K}と小さく、スケートの刃が滑る理由をこれだけで説明することはできません（氷の表面の薄い水の層のでき方は、今も研究されています）。傾きの大きさは、熱-10 のクラペイロン・クラウジウスの式で計算できます。
\end{seikaku}

\bksubsection{2}{わいている途中の水：湿り蒸気と乾き度}

\leavevmode\label{ph:熱-8:sec2}圧力を一定に保ったまま、水\qty{1}{kg}をゆっくり温めていきます。容器はピストンのついた筒で、ピストンには一定の重りがのっていて、中の圧力はいつも同じです（図\ref{ph:熱-8:fig-tsutsu}）。
\begin{bkfigure}
\begin{tikzpicture}[x=1mm, y=1mm]
  \foreach \i/\qxname/\qxtemp in {0/圧縮液/\qty{20}{\celsius},1/飽和液/\qty{99.97}{\celsius},2/湿り蒸気/\qty{99.97}{\celsius},3/飽和蒸気/\qty{99.97}{\celsius},4/過熱蒸気/\qty{150}{\celsius}} {
    \begin{scope}[shift={(\i*26,0)}]
      \draw[bkInk, line width=0.8pt] (0,40) -- (0,0) -- (16,0) -- (16,40);
      \ifnum\i=0 \fill[bkBlue!25] (0.4,0.4) rectangle (15.6,7.6); \filldraw[fill=bkGray!45, draw=bkInk] (0.3,7.6) rectangle (15.7,9.8); \fi
      \ifnum\i=1 \fill[bkBlue!25] (0.4,0.4) rectangle (15.6,7.8); \filldraw[fill=bkGray!45, draw=bkInk] (0.3,7.8) rectangle (15.7,10); \fi
      \ifnum\i=2 \fill[bkBlue!25] (0.4,0.4) rectangle (15.6,4); \fill[bkBlue!7] (0.4,4) rectangle (15.6,20); \foreach \x in {4,8,12} {\draw[bkBlue!70!black, line width=0.4pt] (\x,2.5) circle (0.7);} \filldraw[fill=bkGray!45, draw=bkInk] (0.3,20) rectangle (15.7,22.2); \fi
      \ifnum\i=3 \fill[bkBlue!7] (0.4,0.4) rectangle (15.6,31); \filldraw[fill=bkGray!45, draw=bkInk] (0.3,31) rectangle (15.7,33.2); \fi
      \ifnum\i=4 \fill[bkRed!7] (0.4,0.4) rectangle (15.6,36); \filldraw[fill=bkGray!45, draw=bkInk] (0.3,36) rectangle (15.7,38.2); \fi
      \node[qxi lab, anchor=north] at (8,-1.5) {\qxname};
      \node[qxi lab, anchor=north] at (8,-6.2) {\qxtemp};
      \draw[qxi heat, line width=1pt] (8,-16) -- (8,-11);
    \end{scope}}
  \node[qxi lab, anchor=west] at (122,-13.5) {熱};
\end{tikzpicture}
\bkfigcaption{1気圧（\qty{101.325}{kPa}）で、水\qty{1}{kg}を下から温めていく（ピストンの重さで圧力を一定に保つ）。体積の目もりは正しくない（本当の水蒸気の体積は、水の約1600倍）。}
\label{ph:熱-8:fig-tsutsu}
\end{bkfigure}
\begin{enumerate}[label={\bkcircled[bkBlue]{\arabic*}}, leftmargin=2.4\zw, itemsep=2pt, topsep=2pt]
\item \textbf{圧縮液}\bkidx{あっしゅくえき}{圧縮液}：まだわく温度に届いていない水。温めると温度が上がる。
\item \textbf{飽和液}\bkidx{ほうわえき}{飽和液}：ちょうどわく温度（1気圧なら\qty{99.97}{\celsius}）になった水。
\item \textbf{湿り蒸気}\bkidx{しめりじょうき}{湿り蒸気}：わいている途中。水と水蒸気が混ざっていて、温度は一定。
\item \textbf{飽和蒸気}\bkidx{ほうわじょうき}{飽和蒸気}：水が全部、水蒸気になった瞬間。温度はまだわく温度。
\item \textbf{過熱蒸気}：さらに温めると、水蒸気の温度が上がる。
\end{enumerate}
ここで「圧縮液」は、わく温度より冷たい水のことで、強く押し縮めた水という意味ではありません。図\ref{ph:熱-8:fig-kanetsu}は、加えた熱と温度の関係です。熱-1 の、氷を温めるグラフと同じ形をしていて、わいている間（③）は温度が平らです。
\begin{bkfigure}
\begin{tikzpicture}
\begin{axis}[bkplot, width=100mm, height=52mm, xmin=0, xmax=2900, ymin=0, ymax=185, xtick={0,500,1000,1500,2000,2500}, ytick={20,50,100,150}, xlabel={熱（kJ）}, ylabel={温度（℃）}]
  \addplot[bkcurve] coordinates {(0,20) (335,99.97) (2592,99.97) (2693,150)};
  \addplot[only marks, mark=*, mark size=1.6pt, bkOrange] coordinates {(335,99.97) (2592,99.97)};
  \draw[bkGray, line width=0.5pt, {Stealth[length=3pt]}-{Stealth[length=3pt]}] (axis cs:335,120) -- (axis cs:2592,120);
  \node[font=\sffamily\footnotesize, anchor=south] at (axis cs:1463,122) {わいている間：\qty{2256.5}{kJ}（蒸発熱 $L$）};
  \node[font=\sffamily\footnotesize, anchor=north west] at (axis cs:200,58) {水を温める：\qty{335}{kJ}};
  \node[font=\sffamily\footnotesize, anchor=east] at (axis cs:2610,160) {水蒸気を温める：\qty{101}{kJ}};
  \node[font=\sffamily\footnotesize, anchor=north west] at (axis cs:380,95) {飽和液};
  \node[font=\sffamily\footnotesize, anchor=north] at (axis cs:2500,95) {飽和蒸気};
\end{axis}
\end{tikzpicture}
\bkfigcaption{1気圧で\qty{20}{\celsius}の水\qty{1}{kg}を\qty{150}{\celsius}の水蒸気にするとき、加えた熱（横）と温度（縦）。わかすのに使う熱が、ほとんどをしめる。}
\label{ph:熱-8:fig-kanetsu}
\end{bkfigure}

\noindent\textbf{乾き度：何割が水蒸気か}\par\nobreak
湿り蒸気\qty{1}{kg}のうち、水蒸気になった部分の質量の割合を\textbf{乾き度} $x$ といいます。\bkidx{かわきど}{乾き度}飽和液は $x=0$、飽和蒸気は $x=1$ で、その間は0と1の間の数です（2-5③の割合）。たとえば $x=0.9$ は、「1 kg のうち\qty{0.9}{kg}が水蒸気、\qty{0.1}{kg}が水」という意味です。

蒸気表では、飽和液の量に「$'$」（ダッシュ）、飽和蒸気の量に「$''$」（ツーダッシュ）をつけて書きます。たとえば $h'$ は飽和液の比エンタルピー、$h''$ は飽和蒸気の比エンタルピーです（比エンタルピー $h$ は熱-4 で学んだ「1 kg あたりのエンタルピー」）。湿り蒸気 1 kg のうち、$(1-x)$〔kg〕は $h'$ の水、$x$〔kg〕は $h''$ の水蒸気なので、合わせると
\[ h=(1-x)h'+xh'' \]
です。かっこをはずして、$h'$ をまとめると
\pheq{h=h'+x\,(h''-h')}{kJ/kg ＝ kJ/kg}
となります。\label{ph:熱-8:kawaki}比体積 $v$ や比エントロピー $s$ も同じ形です。
\[ v=v'+x\,(v''-v'),\qquad s=s'+x\,(s''-s') \]
\begin{phyomi}[$h=h'+x(h''-h')$]
「飽和液の $h'$ から出発して、わかした割合 $x$ の分だけ、蒸発熱 $h''-h'$ を足す」。2つの値 $h'$ と $h''$ の、重みつきの平均です。$x=0$ で $h'$、$x=1$ で $h''$ になります。
\end{phyomi}
\begin{phbaisuru}{乾き度 $x$}
$h$ は2倍にはなりません。$h$ は $x$ の1次関数（3-5）で、$x$ が増えた分だけ、蒸発熱 $h''-h'$ の分が比例して増えるからです。\qty{10}{kPa}の湿り蒸気（$h'=\qty{191.8}{kJ/kg}$、$h''=\qty{2583.9}{kJ/kg}$）では、$x=0.4$ で\qty{1148.6}{kJ/kg}、$x=0.8$ で\qty{2105.5}{kJ/kg}です。増えたのは「蒸発した分」の\qty{956.9}{kJ/kg}で、2倍の\qty{2297.2}{kJ/kg}にはなりません。
\par\vspace{2pt}\centering
\begin{tikzpicture}
\begin{axis}[bkplot, width=96mm, height=33mm, xmin=0, xmax=1.05, ymin=0, ymax=2900, xtick={0,0.2,...,1}, ytick={0,1000,2000}, xlabel={乾き度 $x$}, ylabel={$h$（kJ/kg）}]
  \addplot[bkcurve, domain=0:1] {191.8+x*(2583.9-191.8)};
  \addplot[only marks, mark=*, mark size=1.6pt, bkOrange] coordinates {(0,191.8) (0.4,1148.6) (0.8,2105.5) (1,2583.9)};
  \node[font=\sffamily\footnotesize, anchor=north west] at (axis cs:0.41,1120) {$x=0.4$：1148.6};
  \node[font=\sffamily\footnotesize, anchor=north west] at (axis cs:0.81,2080) {$x=0.8$：2105.5};
  \node[font=\sffamily\footnotesize, anchor=south west] at (axis cs:0.02,260) {$h'$};
  \node[font=\sffamily\footnotesize, anchor=south east] at (axis cs:0.98,2620) {$h''$};
\end{axis}
\end{tikzpicture}
\end{phbaisuru}

\noindent\textbf{蒸発熱は $h''-h'$}\par\nobreak
飽和液 1 kg を、同じ圧力のまま飽和蒸気にするのに要る熱は、熱-4 のバランス式で求められます。わかす装置に、1 kg あたり熱 $q_{\text{入}}$ が入り、仕事の出入りはなく、入口と出口で速さ・高さの差がないとすると
\[ q_{\text{入}}-0=h''-h' \]
です。この熱が、熱-1 の蒸発熱 $L$ です。熱-1 では、加えた熱を測って $L$ を決めました。蒸気表では、$L$ が $h''$ と $h'$ の差として読めます。
\Needspace{13\baselineskip}
\begin{phflow}
\phflowrow{$q_{\text{入}}-w_{\text{出}}=h_2-h_1$\\（熱-4 のバランス式）}{仕事なし・圧力一定\\飽和液→飽和蒸気}{$L=h''-h'$}
\phflowrow{$L=h''-h'$}{質量 $m$ を\\全部わかす}{$Q=mL$\\（熱-1）}
\phflowrow{$q_{\text{入}}-w_{\text{出}}=h_2-h_1$}{仕事なし\\飽和液→乾き度 $x$}{$q_{\text{入}}=xL$}
\end{phflow}
表\ref{ph:熱-8:tab-ps}の $L$ は、温度が高いほど小さくなり、臨界点で0になります。臨界点では、水と水蒸気の区別がなくなるからです。

\bksubsection{3}{蒸気表の読み方と、表のあいだの値}

\leavevmode\label{ph:熱-8:sec3}発電所では、温度よりも圧力で状態を決めることが多いので、圧力ごとに並べた飽和表も使います（表\ref{ph:熱-8:tab-sat}）。\bkidx{じょうきひょう}{蒸気表}表\ref{ph:熱-8:tab-ps}と同じ IAPWS-IF97（国際水・蒸気性質協会が1997年に決めた、工業用の水と水蒸気の式）から計算した値で、巻末の蒸気表の抜き出しです。
\begin{bkcenter}
\begin{bktable}{colspec={Q[c,12mm] Q[r,11mm] Q[r,15mm] Q[r,13mm] Q[r,12mm] Q[r,12mm] Q[r,18mm] Q[r,18mm]}, colsep=2.6pt, rowsep=1.2pt, row{1}={font=\sffamily\footnotesize\bfseries, halign=c}, cells={font=\footnotesize}}
{$p$\\〔MPa〕} & {$\theta$\\〔℃〕} & {$v'$\\〔m$^3$/kg〕} & {$v''$\\〔m$^3$/kg〕} & {$h'$\\〔kJ/kg〕} & {$h''$\\〔kJ/kg〕} & {$s'$\\〔kJ/(kg·K)〕} & {$s''$\\〔kJ/(kg·K)〕} \\
0.01 & 45.81 & 0.001010 & 14.67 & 191.8 & 2583.9 & 0.6492 & 8.1489 \\
0.1 & 99.61 & 0.001043 & 1.694 & 417.4 & 2674.9 & 1.3026 & 7.3588 \\
0.2 & 120.21 & 0.001061 & 0.8857 & 504.7 & 2706.2 & 1.5301 & 7.1269 \\
0.5 & 151.84 & 0.001093 & 0.3748 & 640.2 & 2748.1 & 1.8606 & 6.8206 \\
1 & 179.89 & 0.001127 & 0.1943 & 762.7 & 2777.1 & 2.1384 & 6.5850 \\
2 & 212.38 & 0.001177 & 0.09958 & 908.6 & 2798.4 & 2.4470 & 6.3392 \\
4 & 250.36 & 0.001253 & 0.04978 & 1087.4 & 2800.9 & 2.7967 & 6.0697 \\
10 & 311.00 & 0.001453 & 0.01803 & 1407.9 & 2725.5 & 3.3603 & 5.6159 \\
\end{bktable}
\bktabcaption{水の飽和表（圧力基準。IAPWS-IF97 の式から計算）。\qty{0.1}{MPa}は\qty{100}{kPa}で、1気圧（\qty{101.325}{kPa}）より少し低い。}
\label{ph:熱-8:tab-sat}
\end{bkcenter}
\noindent 表の読み方を、\qty{1}{MPa}の行でたしかめましょう。
\begin{enumerate}[label={\bkcircled[bkBlue]{\arabic*}}, leftmargin=2.4\zw, itemsep=2pt, topsep=2pt]
\item $\theta=\qty{179.89}{\celsius}$：\qty{1}{MPa}（約10気圧）では、水は約\qty{180}{\celsius}でわく。
\item $v'=\qty{0.001127}{m^3/kg}$、$v''=\qty{0.1943}{m^3/kg}$：わくと、体積が約170倍になる。
\item $h'=\qty{762.7}{kJ/kg}$、$h''=\qty{2777.1}{kJ/kg}$：蒸発熱は $L=h''-h'=\qty{2014.4}{kJ/kg}$。
\item $s'$、$s''$：熱-6 の比エントロピー。⑤⑥でタービンの計算に使う。
\end{enumerate}
$h$ と $s$ の0の点は、三重点の水（\qty{0.01}{\celsius}の飽和液）で $h'\approx0$、$s'=0$ になるように決めてあります（国際的な約束）。だから、$h$ の値そのものより、2つの状態の\textbf{差}に意味があります。

\noindent\textbf{表にない圧力：直線補間}\par\nobreak
表にない\qty{150}{kPa}の飽和温度がほしいときは、3-5⑧のように、となりの2行を直線で結んで読みます（\textbf{直線補間}）。\bkidx{ちょくせんほかん}{直線補間}2行を $(x_1,\ y_1)$、$(x_2,\ y_2)$ とすると、その間の $x$ での値は
\pheq{y=y_1+(y_2-y_1)\times\frac{x-x_1}{x_2-x_1}}{y の単位 ＝ y の単位}
です。$x-x_1$ を $x_2-x_1$ でわった値は「$x_1$ から $x_2$ までの何割まで来たか」で、乾き度の式と同じ「重みつきの平均」の形です。\bkreidairef{3}では、この読み方の誤差がどれくらいかも確かめます。
\begin{bkcenter}
\bksheet{4}{24}{%
1 & $p$〔kPa〕 & $\theta$〔℃〕 & 割合 & 補間の $\theta$〔℃〕 \\
2 & 100 & 99.61 &  &  \\
3 & 200 & 120.21 &  &  \\
4 & 150 &  & 0.5 & 109.91 \\}
\bktabcaption{表計算で補間する。C4 は \texttt{=(A4-A2)/(A3-A2)}、D4 は \texttt{=B2+(B3-B2)*C4}。}
\end{bkcenter}

\begin{seikaku}
直線補間は、2つの点の間の線が「ほぼまっすぐ」なときに使える読み方です。飽和蒸気圧のように、温度で急に（倍々に近く）増える量を、温度で補間すると大きくずれます。\qty{20}{\celsius}（\qty{2.339}{kPa}）と\qty{40}{\celsius}（\qty{7.384}{kPa}）から\qty{30}{\celsius}を直線で補間すると\qty{4.862}{kPa}ですが、本当は\qty{4.247}{kPa}で、14\,\% も大きくなります。片対数グラフ（3-12⑥）で直線になる量は、対数で補間する（$\sqrt{2.339\times7.384}\approx\qty{4.156}{kPa}$）か、表の間かくをもっと細かくします。本物の蒸気表が、温度や圧力を細かく刻んでいるのはこのためです。
\end{seikaku}

\bksubsection{4}{わいたあとの水蒸気：過熱蒸気の表}

\leavevmode\label{ph:熱-8:sec4}飽和蒸気をさらに温めた過熱蒸気は、圧力と温度の2つを決めれば状態が決まります。そこで、過熱蒸気表は、圧力ごとに、温度を縦に並べた表になっています（表\ref{ph:熱-8:tab-kanetsu}）。\bkidx{かねつじょうき}{過熱蒸気}飽和温度より何 K 高いかを\textbf{過熱度}といいます。たとえば\qty{1}{MPa}・\qty{300}{\celsius}の蒸気の過熱度は、$300-179.89\approx\qty{120}{K}$ です。
\Needspace{7\baselineskip}
\begin{phmijika}[工場の蒸気・スチームアイロン]
食品工場や病院では、ボイラーでつくった\qtyrange{0.5}{1}{MPa}くらいの蒸気を管で送り、温めたり、器具を消毒したりしています。蒸気が管の中で冷えて水にもどる（凝縮する）とき、蒸発熱 $h''-h'$ の大きな熱を出すので、少しの量で多くの熱を運べるからです。スチームアイロンが布のしわをすばやく伸ばせるのも、布にふれた水蒸気が水にもどるときに、布に大きな熱をわたすからです。
\end{phmijika}

\begin{seikaku}
水蒸気も気体なので、理想気体の式（熱-3）$pv=R_{\mathrm s}T$（水蒸気の $R_{\mathrm s}=\qty{461.5}{J/(kg.K)}$）で比体積を計算してみます。\qty{1}{MPa}・\qty{300}{\celsius}では $461.5\times573.15\div(1\times10^6)\approx\qty{0.2645}{m^3/kg}$ で、表\ref{ph:熱-8:tab-kanetsu}の\qty{0.2580}{m^3/kg}より2.5\,\% 大きくなります。わく温度に近い所や、圧力が高い所では、分子どうしが引き合う力がきいて、理想気体からずれるのです。\qty{10}{kPa}・\qty{100}{\celsius}のように圧力が低ければ、ずれは0.1\,\% くらいです。蒸気の計算に、理想気体の式ではなく蒸気表を使うのはこのためです。
\end{seikaku}

\begin{bkcenter}
\begin{bktable}{colspec={Q[c,12mm] Q[c,11mm] Q[r,16mm] Q[r,15mm] Q[r,19mm]}, colsep=4pt, rowsep=1.2pt, row{1}={font=\sffamily\footnotesize\bfseries, halign=c}, cells={font=\footnotesize}}
{$p$\\〔MPa〕} & {$\theta$\\〔℃〕} & {$v$\\〔m$^3$/kg〕} & {$h$\\〔kJ/kg〕} & {$s$\\〔kJ/(kg·K)〕} \\
1 & 200 & 0.2060 & 2828.3 & 6.6955 \\
1 & 250 & 0.2327 & 2943.2 & 6.9266 \\
1 & 300 & 0.2580 & 3051.7 & 7.1247 \\
1 & 350 & 0.2825 & 3158.2 & 7.3028 \\
1 & 400 & 0.3066 & 3264.4 & 7.4668 \\
1 & 500 & 0.3541 & 3479.0 & 7.7640 \\
4 & 300 & 0.05887 & 2961.7 & 6.3638 \\
4 & 400 & 0.07343 & 3214.4 & 6.7712 \\
4 & 500 & 0.08644 & 3445.8 & 7.0919 \\
4 & 600 & 0.09886 & 3674.8 & 7.3704 \\
10 & 350 & 0.02244 & 2924.0 & 5.9458 \\
10 & 400 & 0.02644 & 3097.4 & 6.2139 \\
10 & 500 & 0.03281 & 3375.1 & 6.5993 \\
10 & 600 & 0.03838 & 3625.8 & 6.9045 \\
\end{bktable}
\bktabcaption{過熱蒸気表の抜き出し（IAPWS-IF97 の式から計算）。}
\label{ph:熱-8:tab-kanetsu}
\end{bkcenter}

\bksubsection{5}{状態を図で見る：$T$-$s$ 線図と $h$-$s$ 線図}

\leavevmode\label{ph:熱-8:sec5}熱-6 の $T$-$s$ 線図に、表\ref{ph:熱-8:tab-ps}・表\ref{ph:熱-8:tab-sat}の飽和液（$s'$）と飽和蒸気（$s''$）の点を、いろいろな温度でかきこむと、山の形の線ができます（図\ref{ph:熱-8:fig-ts}）。この山を\textbf{飽和ドーム}とよびます。山の左のすそが飽和液の線、右のすそが飽和蒸気の線、頂上が臨界点です。山の中は湿り蒸気、左の外は圧縮液、右の外は過熱蒸気です。
\begin{bkfigure}
\begin{tikzpicture}
\begin{axis}[bkplot, width=104mm, height=72mm, xmin=0, xmax=9.4, ymin=0, ymax=560, xtick={0,1,...,9}, ytick={0,100,...,500}, xlabel={$s$（kJ/(kg·K)）}, ylabel={温度（℃）}]
  \addplot[bkcurve, smooth] coordinates {(0,0.01) (0.2965,20) (0.6493,45.81) (1.3070,100) (1.8420,150) (2.3308,200) (2.7934,250) (3.2547,300) (3.6599,340) (3.9164,360) (4.1142,370) (4.2380,373) (4.4120,373.95)};
  \addplot[bkcurve, smooth] coordinates {(4.4120,373.95) (4.6293,373) (4.7996,370) (5.0527,360) (5.3359,340) (5.7058,300) (6.0722,250) (6.4303,200) (6.8370,150) (7.3541,100) (8.1488,45.81) (8.6661,20) (9.1555,0.01)};
  \addplot[bkcurve2, line width=2pt, smooth] coordinates {(0.6493,46.1) (1.2994,100) (1.8315,150) (2.3177,200) (2.7791,250) (3.2484,300) (3.3603,311.0)};
  \addplot[bkcurve2, line width=2pt] coordinates {(3.3603,311.0) (5.6160,311.0)};
  \addplot[bkcurve2, line width=2pt, smooth] coordinates {(5.6160,311.0) (5.9458,350) (6.2139,400) (6.4217,450) (6.5993,500)};
  \addplot[bkcurve2, line width=2pt] coordinates {(6.5993,500) (6.5993,45.81) (0.6493,45.81)};
  \addplot[only marks, mark=*, mark size=1.8pt, bkInk] coordinates {(6.5993,500) (6.5993,45.81) (0.6493,45.81) (4.4120,373.95)};
  \node[font=\sffamily\footnotesize\bfseries, anchor=west] at (axis cs:6.68,500) {1};
  \node[font=\sffamily\footnotesize\bfseries, anchor=north west] at (axis cs:6.62,44) {2};
  \node[font=\sffamily\footnotesize\bfseries, anchor=north west] at (axis cs:0.75,40) {3・4};
  \node[font=\sffamily\footnotesize, anchor=south] at (axis cs:4.41,385) {臨界点};
  \node[font=\sffamily\footnotesize, align=center] at (axis cs:4.5,118) {湿り蒸気\\（ドームの中）};
  \node[font=\sffamily\footnotesize, anchor=east] at (axis cs:2.4,330) {飽和液の線};
  \node[font=\sffamily\footnotesize, anchor=west] at (axis cs:7.5,170) {飽和蒸気の線};
  \node[font=\sffamily\footnotesize, anchor=west] at (axis cs:7.1,420) {過熱蒸気};
  \node[font=\sffamily\footnotesize, anchor=north] at (axis cs:4.45,232) {\qty{10}{MPa}・\qty{311}{\celsius}でわく};
\end{axis}
\end{tikzpicture}
\bkfigcaption{水の $T$-$s$ 線図（横は比エントロピー $s$、縦は温度。目もりは℃）と飽和ドーム（実線）。太い線は⑥のランキンサイクル（\qty{10}{MPa}・\qty{500}{\celsius}→\qty{10}{kPa}）。番号は図\ref{ph:熱-8:fig-sochi}の装置の番号と同じ。}
\label{ph:熱-8:fig-ts}
\end{bkfigure}
ドームの中では、等しい圧力の線が水平（温度一定）になります。わいている間は温度が変わらないことが、図でも見えます。そして、熱-6 で学んだように、$T$-$s$ 線図で可逆な道の下の面積は、受けとった熱です。わいている間の熱は、高さ $T$〔K〕、はば $s''-s'$ の長方形の面積で、$L=T(s''-s')$ です。\qty{100}{\celsius}で確かめると $373.15\times(7.3541-1.3070)\approx\qty{2256.5}{kJ/kg}$ で、表\ref{ph:熱-8:tab-ps}の $L$ と同じです。

\noindent\textbf{$h$-$s$ 線図（モリエ線図）}\par\nobreak
タービンの計算では、縦に $h$、横に $s$ をとった線図が便利です（図\ref{ph:熱-8:fig-hs}）。ドイツのモリエが1904年に発表したので、\textbf{モリエ線図}ともいいます。\bkidx{もりえせんず}{モリエ線図}タービンの中で蒸気が断熱・可逆に膨張すると、$s$ は一定（熱-6）なので、状態は線図の上を\textbf{真下}へ進みます。タービンが取り出す仕事は $h$ の差（熱-4）なので、その線の\textbf{長さ}が、そのまま仕事になります。
\begin{bkfigure}
\begin{tikzpicture}
\begin{axis}[bkplot, axis lines=left, width=104mm, height=76mm, xmin=5.5, xmax=8.7, ymin=1800, ymax=3800, xtick={6,6.5,...,8.5}, ytick={2000,2400,2800,3200,3600}, xlabel={$s$（kJ/(kg·K)）}, ylabel={$h$（kJ/kg）}]
  \addplot[bkGray, line width=0.8pt, smooth] coordinates {(5.899,1866) (6.649,2105) (7.399,2345) (8.149,2584) (8.321,2641) (8.576,2737)};
  \addplot[bkGray, line width=0.8pt, smooth] coordinates {(5.542,1998) (6.148,2223) (6.753,2449) (7.359,2675) (7.516,2736) (7.749,2835) (8.033,2974) (8.545,3279)};
  \addplot[bkGray, line width=0.8pt, smooth] coordinates {(5.696,2374) (6.140,2576) (6.585,2777) (6.745,2852) (6.968,2965) (7.233,3115) (7.467,3264) (7.764,3479)};
  \addplot[bkGray, line width=0.8pt, smooth] coordinates {(5.742,2630) (6.070,2801) (6.261,2904) (6.503,3043) (6.771,3214) (7.092,3446) (7.370,3675)};
  \addplot[bkGray, line width=0.8pt, smooth] coordinates {(5.616,2725) (5.885,2886) (6.172,3069) (6.463,3272) (6.599,3375) (6.905,3626)};
  \addplot[bkInk, line width=1.6pt, smooth] coordinates {(5.745,2759) (6.070,2801) (6.339,2798) (6.585,2777) (6.821,2748) (7.359,2675) (7.593,2645) (8.149,2584) (8.394,2561)};
  \draw[bkRed!80!black, line width=1.6pt, -{Stealth[length=4.5pt]}] (axis cs:6.5993,3375.1) -- (axis cs:6.5993,2098);
  \addplot[only marks, mark=*, mark size=2pt, bkRed!80!black] coordinates {(6.5993,3375.1) (6.5993,2089.6)};
  \node[qxi lab, anchor=east] at (axis cs:6.56,3400) {1};
  \node[qxi lab, anchor=north west] at (axis cs:6.62,2080) {2};
  \node[qxi lab, anchor=south east] at (axis cs:6.90,3640) {\qty{10}{MPa}};
  \node[qxi lab, anchor=south east] at (axis cs:7.37,3690) {\qty{4}{MPa}};
  \node[qxi lab, anchor=south east] at (axis cs:7.76,3490) {\qty{1}{MPa}};
  \node[qxi lab, anchor=south east] at (axis cs:8.54,3290) {\qty{0.1}{MPa}};
  \node[qxi lab, anchor=south east] at (axis cs:8.57,2750) {\qty{10}{kPa}};
  \node[font=\sffamily\small\bfseries] at (axis cs:8.28,3600) {過熱蒸気};
  \node[font=\sffamily\small\bfseries] at (axis cs:7.65,2050) {湿り蒸気};
\end{axis}
\end{tikzpicture}
\bkfigcaption{水蒸気の $h$-$s$ 線図（モリエ線図。横は比エントロピー $s$、縦は比エンタルピー $h$）。細い灰色の線は等しい圧力の線、黒い太い線が飽和蒸気の線（$x=1$）で、その下が湿り蒸気、上が過熱蒸気。矢印は\bkreidairef{6}のタービン（1：\qty{10}{MPa}・\qty{500}{\celsius} → 2：\qty{10}{kPa}・$x=0.79$）で、長さが仕事\qty{1285}{kJ/kg}。}
\label{ph:熱-8:fig-hs}
\end{bkfigure}

\bksubsection{6}{水蒸気で回す：ランキンサイクル}

\leavevmode\label{ph:熱-8:sec6}火力発電所・原子力発電所・ごみ焼却場の発電機は、水を水蒸気にしてタービンを回し、その水蒸気を水にもどして、また使います。いちばん簡単な形が、4つの機械をつないだ\textbf{ランキンサイクル}です（図\ref{ph:熱-8:fig-sochi}）。\bkidx{らんきんさいくる}{ランキンサイクル}スコットランドの技術者ランキンが、1859年の本で、蒸気機関の計算のしかたをまとめました。
\begin{itemize}[label={\color{bkTheme}●}, leftmargin=1.6\zw, itemsep=2pt, topsep=2pt]
\item \textbf{4→1 ボイラー}：高い圧力の水を温めて、わかし、過熱蒸気にする。熱 $q_{\text{入}}$ を受けとる。
\item \textbf{1→2 タービン}：蒸気が羽根車を回しながら膨張し、圧力と温度が下がる。仕事 $w_{\text{出}}$ を取り出す。
\item \textbf{2→3 復水器}：タービンを出た湿り蒸気を、海や川の水で冷やして水にもどす。熱 $q_{\text{出}}$ を捨てる。
\item \textbf{3→4 ポンプ}：水を押して、ボイラーの圧力まで上げる。仕事 $w_{\text{入}}$ を入れる。
\end{itemize}
この章では、熱と仕事をどれも\textbf{正の数}で書き、向きを「入」「出」の添字で表します（熱-5 の $Q_{\mathrm H}$・$Q_{\mathrm L}$ と同じ書き方）。
\begin{bkfigure}
\qxirankine
\bkfigcaption{ランキンサイクルの装置。管の中の矢印が水・水蒸気の流れ。外から出入りする矢印のうち、$q$ の札が熱、$w$ の札が仕事。丸の番号は状態の番号。}
\label{ph:熱-8:fig-sochi}
\end{bkfigure}

\noindent\textbf{それぞれの機械にバランス式を使う}\par\nobreak
4つの機械は、どれも水・水蒸気が流れぬける機械です。熱-4 の定常流動のエネルギーの式を、入口と出口で速さ・高さの差を無視して、それぞれに使います。
\begin{phmodel}[理想のランキンサイクル]
\phignore{管の中の圧力の下がり（摩擦）。機械から周りへ逃げる熱。入口と出口の速さ・高さの差。}
\phassume{流れは定常。タービンとポンプは断熱で可逆（$s$ が一定）。ボイラーと復水器の中は圧力が一定。復水器の出口は飽和液。}
\phlimit{実際のタービンの中の摩擦・渦（$s$ が増える：⑧）。湿り蒸気の水の粒がタービンの羽根に当たって失う分。}
\end{phmodel}
\begin{enumerate}[label={\bkcircled[bkBlue]{\arabic*}}, leftmargin=2.4\zw, itemsep=2pt, topsep=2pt]
\item ボイラー（仕事なし）：$q_{\text{入}}=h_1-h_4$
\item タービン（熱の出入りなし）：$w_{\text{出}}=h_1-h_2$
\item 復水器（仕事なし）：$q_{\text{出}}=h_2-h_3$
\item ポンプ（熱の出入りなし）：$w_{\text{入}}=h_4-h_3$
\end{enumerate}
1 kg の水が4つの機械を1周してもとにもどると、比エンタルピーももとにもどるので、入ったエネルギーと出たエネルギーが等しくなります。
\[ q_{\text{入}}+w_{\text{入}}=w_{\text{出}}+q_{\text{出}} \]
取り出せる正味の仕事は $w_{\text{出}}-w_{\text{入}}=q_{\text{入}}-q_{\text{出}}$ で（熱-5 の $W_{\text{出}}$ にあたる量。この章の $w_{\text{出}}$ はタービンの仕事）、効率（理論熱効率）は
\pheq{\eta=\frac{w_{\text{出}}-w_{\text{入}}}{q_{\text{入}}}=1-\frac{q_{\text{出}}}{q_{\text{入}}}}{(kJ/kg)/(kJ/kg)：単位なし}
です。\label{ph:熱-8:eta}熱-5 の熱機関の効率 $\eta=1-Q_{\mathrm L}/Q_{\mathrm H}$ と同じ形で、1 kg あたりで書いたものです。

\noindent\textbf{ポンプの仕事は $v\,\Delta p$}\par\nobreak
液体の水は、押してもほとんど縮みません。熱-4 で学んだように、流れる物の比エンタルピーには「流れの仕事」$pv$ の分（$v=1/\rho$）がふくまれています。縮まない水を断熱・可逆に押すときは、温度も内部エネルギーもほとんど変わらず、増えるのはこの $pv$ の分だけなので
\pheq{w_{\text{入}}=h_4-h_3\approx v_3\,(p_4-p_3)}{kJ/kg ＝ m$^3$/kg × kPa}
となります。\label{ph:熱-8:pump}力-13 のポンプの「$\Delta p/\rho$」と同じ式です。水の $v$ は約\qty{0.001}{m^3/kg}と小さいので、\qty{10}{MPa}まで押しても1 kg あたり約\qty{10}{kJ}です（\bkreidairef{7}）。同じ圧力比で気体を押す圧縮機（熱-4）よりずっと小さく、タービンの仕事の1\,\% にもなりません。いったん水にもどしてから押すのは、このためです。
\Needspace{12\baselineskip}
\begin{phflow}
\phflowrow{$q_{\text{入}}-w_{\text{出}}=h_2-h_1$\\（熱-4）}{機械ごとに\\仕事なし・断熱}{$q_{\text{入}}=h_1-h_4$\quad$w_{\text{出}}=h_1-h_2$\\$q_{\text{出}}=h_2-h_3$\quad$w_{\text{入}}=h_4-h_3$}
\phflowrow{$w_{\text{入}}=h_4-h_3$}{水は縮まない\\断熱・可逆}{$w_{\text{入}}\approx v_3(p_4-p_3)$}
\end{phflow}

\noindent\textbf{タービンの出口の状態：$s$ が同じ点を探す}\par\nobreak
タービンの出口2は、圧力 $p_2$（復水器の圧力）と、$s_2=s_1$（断熱・可逆）で決まります。2はふつうドームの中（湿り蒸気）なので、②の $s=s'+x(s''-s')$ を $x$ について解いて
\pheq{x_2=\frac{s_1-s'}{s''-s'}}{(kJ/(kg·K))/(kJ/(kg·K))：単位なし}
で乾き度を出し、$h_2=h'+x_2(h''-h')$ を計算します（$s'$・$s''$・$h'$・$h''$ は、圧力 $p_2$ の飽和表の値）。\bkreidairef{6}で、\qty{10}{MPa}・\qty{500}{\celsius}の蒸気を\qty{10}{kPa}まで膨張させるサイクルを、この手順で計算すると、効率は約40\,\% になります。

\noindent\textbf{カルノーの効率より低いのはなぜ}\par\nobreak
このサイクルの最高の温度は\qty{500}{\celsius}（\qty{773}{K}）、最低は復水器の\qty{45.8}{\celsius}（\qty{319}{K}）です。熱-5 のカルノーの効率は $1-319/773\approx0.59$ で、ランキンの0.40より大きくなります。図\ref{ph:熱-8:fig-ts}を見ると、わけがわかります。ボイラーでは、水が\qty{46}{\celsius}から少しずつ温まり、\qty{311}{\celsius}でわいて、最後に\qty{500}{\celsius}になります。熱を受けとる温度の多くが、\qty{500}{\celsius}よりずっと低いのです。

受けとった熱 $q_{\text{入}}$ を、受けとる間の $s$ の増え $s_1-s_4$ でわった値を、\textbf{平均の受熱温度} $T_{\mathrm m}$ とよびます（$T$-$s$ 線図で、受熱の道の下の面積と同じ面積の長方形の高さ）。
\[ T_{\mathrm m}=\frac{q_{\text{入}}}{s_1-s_4}=\frac{3173.2}{6.5993-0.6492}\approx\qty{533}{K}\quad(\qty{260}{\celsius}) \]
熱を捨てるのは $T_{\mathrm L}=\qty{319}{K}$ の一定の温度なので、$\eta=1-T_{\mathrm L}/T_{\mathrm m}=1-319/533\approx0.40$ と、ランキンの効率にぴったり合います。効率を上げるには、$T_{\mathrm m}$ を上げる（高い温度で熱を受けとる）か、$T_{\mathrm L}$ を下げればよいのです。⑦の再熱・再生と、⑧の超臨界圧は、どれも $T_{\mathrm m}$ を上げるくふうです。

\bksubsection{7}{効率を上げるくふう：再熱と再生}

\noindent\label{ph:熱-8:sec7}\textbf{再熱：もう一度温める}\par\nobreak
ボイラーの圧力を上げると、わく温度が上がるので $T_{\mathrm m}$ が上がり、効率は上がります。ところが、図\ref{ph:熱-8:fig-hs}でわかるように、圧力が高いほど、同じ温度から真下へ膨張したときの出口が、ドームの深い所（乾き度の小さい所）になります。乾き度0.79は、2割以上が水の粒です。高速で回るタービンの羽根に水の粒が当たると、羽根がけずられ（エロージョン）、効率も下がります。ふつうは、出口の乾き度を約0.88より大きく保つように設計します。

そこで、タービンをとちゅうで2つに分け、途中まで膨張した蒸気をボイラーにもどして、もう一度温めてから、残りのタービンで膨張させます。これを\textbf{再熱}といいます（図\ref{ph:熱-8:fig-saisei}の「再熱」）。\bkidx{さいねつ}{再熱}出口の乾き度が上がり、高い温度で受けとる熱が増えるので、効率も少し上がります（\bkreidairef{9}、図\ref{ph:熱-8:fig-reheat}）。
\begin{bkfigure}
\begin{tikzpicture}
\begin{axis}[bkplot, width=100mm, height=58mm, xmin=0, xmax=9.4, ymin=0, ymax=560, xtick={0,1,...,9}, ytick={0,100,...,500}, xlabel={$s$（kJ/(kg·K)）}, ylabel={温度（℃）}]
  \addplot[bkcurve, smooth] coordinates {(0,0.01) (0.2965,20) (0.6493,45.81) (1.3070,100) (1.8420,150) (2.3308,200) (2.7934,250) (3.2547,300) (3.6599,340) (3.9164,360) (4.1142,370) (4.2380,373) (4.4120,373.95)};
  \addplot[bkcurve, smooth] coordinates {(4.4120,373.95) (4.6293,373) (4.7996,370) (5.0527,360) (5.3359,340) (5.7058,300) (6.0722,250) (6.4303,200) (6.8370,150) (7.3541,100) (8.1488,45.81) (8.6661,20) (9.1555,0.01)};
  \addplot[bkcurve2, line width=2pt, smooth] coordinates {(0.6493,46.1) (1.2994,100) (1.8315,150) (2.3177,200) (2.7791,250) (3.2484,300) (3.3603,311.0)};
  \addplot[bkcurve2, line width=2pt] coordinates {(3.3603,311.0) (5.6160,311.0)};
  \addplot[bkcurve2, line width=2pt, smooth] coordinates {(5.6160,311.0) (5.9458,350) (6.2139,400) (6.4217,450) (6.5993,500)};
  \addplot[bkcurve2, line width=2pt] coordinates {(6.5993,500) (6.5993,353.8)};
  \addplot[bkcurve2, line width=2pt, smooth] coordinates {(6.5993,353.8) (6.7712,400) (6.9383,450) (7.0919,500)};
  \addplot[bkcurve2, line width=2pt] coordinates {(7.0919,500) (7.0919,45.81) (0.6492,45.81)};
  \addplot[bkGray, densely dashed, line width=0.8pt] coordinates {(6.5993,353.8) (6.5993,45.81)};
  \addplot[only marks, mark=*, mark size=1.8pt, bkInk] coordinates {(6.5993,500) (6.5993,353.8) (7.0919,500) (7.0919,45.81) (0.6493,45.81) (6.5993,45.81)};
  \node[font=\sffamily\footnotesize\bfseries, anchor=east] at (axis cs:6.52,505) {1};
  \node[font=\sffamily\footnotesize\bfseries, anchor=east] at (axis cs:6.52,353.8) {a};
  \node[font=\sffamily\footnotesize\bfseries, anchor=west] at (axis cs:7.17,505) {b};
  \node[font=\sffamily\footnotesize\bfseries, anchor=west] at (axis cs:7.17,64) {c};
  \node[font=\sffamily\footnotesize\bfseries, anchor=south east] at (axis cs:6.5,50) {2};
  \node[font=\sffamily\footnotesize\bfseries, anchor=west] at (axis cs:1.0,22) {3・4};
  \node[font=\sffamily\footnotesize, anchor=west] at (axis cs:7.2,420) {再熱（\qty{4}{MPa}）};
  \node[font=\sffamily\footnotesize, align=center] at (axis cs:4.5,150) {湿り蒸気\\（ドームの中）};
\end{axis}
\end{tikzpicture}
\bkfigcaption{\bkreidairef{9}の再熱サイクルの $T$-$s$ 線図。1→a で\qty{4}{MPa}まで膨張し、a→b で\qty{500}{\celsius}まで温め直し、b→c で\qty{10}{kPa}まで膨張する。破線は再熱しないときの膨張（1→2、図\ref{ph:熱-8:fig-ts}）。出口 c は 2 より右（乾き度が大きい所）にくる。}
\label{ph:熱-8:fig-reheat}
\end{bkfigure}

\noindent\textbf{再生：タービンの蒸気で給水を温める}\par\nobreak
ボイラーに入る水（給水）は\qty{46}{\celsius}くらいと冷たく、ボイラーの熱の一部は、この冷たい水を温めるのに使われています。これが $T_{\mathrm m}$ を下げていました。そこで、タービンのとちゅうから蒸気を少し取り出し（\textbf{抽気}）、その蒸気で給水を温めます。これを\textbf{再生}といいます（図\ref{ph:熱-8:fig-saisei}の「再生」）。\bkidx{さいせい}{再生（サイクル）}抽気した蒸気は、タービンで仕事をするはずだった分が減りますが、その蒸発熱は捨てられずに給水にもどるので、全体の効率が上がります（\bkreidairef{10}）。
\begin{bkfigure}
\begin{tikzpicture}[x=1mm, y=1mm]
  % 再熱
  \filldraw[fill=bkRed!8, draw=bkInk, line width=0.8pt, rounded corners=1.2mm] (0,0) rectangle (20,30);
  \node[qxi lab] at (10,24) {ボイラー};
  \draw[bkInk!60, line width=0.5pt] (2,12) -- (18,12);
  \node[qxi lab] at (10,6) {再熱器};
  \filldraw[fill=bkBlue!10, draw=bkInk, line width=0.8pt] (30,26) -- (30,34) -- (40,37) -- (40,23) -- cycle;
  \filldraw[fill=bkBlue!10, draw=bkInk, line width=0.8pt] (46,22) -- (46,34) -- (60,39) -- (60,17) -- cycle;
  \node[qxi lab] at (35,41) {高圧};
  \node[qxi lab] at (53,43) {低圧};
  \draw[qxi flow, line width=1.1pt] (20,26) -- (25,26) -- (25,30) -- (29.4,30);
  \draw[qxi flow, line width=1.1pt] (35,24.4) -- (35,16) -- (24,16) -- (24,8) -- (20.6,8);
  \draw[qxi flow, line width=1.1pt] (20,3) -- (43,3) -- (43,28) -- (45.4,28);
  \draw[qxi flow, line width=1.1pt] (53,18.6) -- (53,10);
  \node[qxi lab, anchor=north] at (53,9.4) {復水器へ};
  \node[qxi lab, anchor=north] at (30,-3) {再熱（タービンを2つに分ける）};
  % 再生
  \begin{scope}[shift={(74,0)}]
  \filldraw[fill=bkRed!8, draw=bkInk, line width=0.8pt, rounded corners=1.2mm] (0,4) rectangle (16,30);
  \node[qxi lab] at (8,17) {ボイラー};
  \filldraw[fill=bkBlue!10, draw=bkInk, line width=0.8pt] (22,26) -- (22,36) -- (44,42) -- (44,20) -- cycle;
  \draw[qxi flow, line width=1.1pt] (16,26) -- (19,26) -- (19,31) -- (21.4,31);
  \draw[qxi flow, line width=1.1pt] (40,21.2) -- (40,12);
  \node[qxi lab, anchor=west] at (40.6,15) {復水器へ};
  \filldraw[fill=bkOrange!12, draw=bkInk, line width=0.8pt, rounded corners=1mm] (21,1) rectangle (35,13);
  \node[qxi lab, align=center] at (28,7) {給水\\加熱器};
  \draw[qxi flow, line width=1.1pt, bkOrange!85!black] (30,24.6) -- (30,13.6);
  \node[qxi lab, anchor=west] at (30.6,19) {抽気};
  \draw[qxi flow, line width=1.1pt] (56,7) -- (35.6,7);
  \node[qxi lab, anchor=south] at (47,7.5) {水（ポンプ）};
  \draw[qxi flow, line width=1.1pt] (20.4,7) -- (18,7) -- (18,10) -- (16.6,10);
  \node[qxi lab, anchor=north] at (28,-3) {再生（抽気で給水を温める）};
  \end{scope}
\end{tikzpicture}
\bkfigcaption{再熱と再生の装置。どちらも、熱を受けとる平均の温度を上げるくふう。}
\label{ph:熱-8:fig-saisei}
\end{bkfigure}

\noindent 図\ref{ph:熱-8:fig-saisei}の「再生」の給水加熱器は、抽気した蒸気と給水を直接混ぜる形（開放形）です。混ぜた出口が、その圧力の飽和液になるように抽気の量を決めます。抽気の割合を $y$（1 kg の蒸気のうち $y$〔kg〕を抽気）とすると、加熱器のバランス式は
\[ y\,h_6+(1-y)\,h_{3'}=h_{\mathrm f} \]
です（$h_6$：抽気の蒸気、$h_{3'}$：ポンプを出た給水、$h_{\mathrm f}$：抽気の圧力の飽和液 $h'$）。熱-4 の混合水栓の計算と同じ形です。実際の大きな発電所では、給水加熱器を7〜8段もつけて、給水を\qty{250}{\celsius}以上まで温めています。

\bksubsection{8}{実際の発電所}

\noindent\label{ph:熱-8:sec8}\textbf{タービンの中の損失}\par\nobreak
実際のタービンでは、蒸気が羽根の間を通るときの摩擦や渦で、$s$ が増えます（熱-6 のエントロピー生成）。断熱・可逆のときの仕事（$h_1-h_{2\mathrm s}$、$2\mathrm s$ は $s$ が一定の出口）に対する、実際の仕事の割合を\textbf{タービンの断熱効率}といい、大きなタービンで0.85〜0.9くらいです。\bkidx{だんねつこうりつ}{断熱効率}
\pheq{w_{\text{出}}=h_1-h_2=0.85\times(h_1-h_{2\mathrm s})}{kJ/kg ＝ kJ/kg}
⑥のサイクルで断熱効率を0.85とすると、$w_{\text{出}}=0.85\times1285.5\approx\qty{1092.7}{kJ/kg}$、出口の比エンタルピーは $h_2=3375.1-1092.7=\qty{2282.4}{kJ/kg}$ です。出口は湿り蒸気のままなので、$x_2=(2282.4-191.8)\div(2583.9-191.8)\approx0.874$ と、少し乾きます（損失の分が熱になって、水の粒を少し蒸発させた）。効率は $(1092.7-10.1)\div3173.2\approx0.34$ に下がります。

\noindent\textbf{復水器の圧力}\par\nobreak
復水器は、\qtyrange{15}{25}{\celsius}くらいの海水で蒸気を冷やします。熱を移すには温度の差がいるので、中の蒸気は海水より\qty{10}{K}ほど高い約\qty{33}{\celsius}で水にもどり、中の圧力はその温度の飽和蒸気圧の約\qty{5}{kPa}（大気圧の約20分の1）と、真空に近くなっています。$T_{\mathrm L}$ をできるだけ下げるためです。冬は海水が冷たいので、同じ発電所でも効率が少し上がります（章末問題）。

\noindent\textbf{超臨界圧の火力発電所}\par\nobreak
最新の石炭火力発電所は、臨界点（\qty{22.064}{MPa}）より高い約\qty{25}{MPa}で水を温め、\qty{600}{\celsius}くらいの蒸気にしています（超臨界圧・超々臨界圧）。臨界点より上では「わく」ことがなく、水は温度が上がるにつれて、なめらかに蒸気へ変わります。$T_{\mathrm m}$ が高くなるので、\qty{25}{MPa}・\qty{600}{\celsius}・再熱（\qty{4}{MPa}で\qty{600}{\celsius}にもどす）・復水器\qty{5}{kPa}の理想のサイクルの効率は約48\,\% になります。タービンなどの損失や、発電所の中で使う電気を差し引くと、実際に送り出す電気は燃料の熱の約42\,\%（高位発熱量で数えた値）です。

\noindent\textbf{原子力発電所}\par\nobreak
沸騰水型の原子力発電所では、原子炉の中で水を直接わかし、約\qty{7}{MPa}・約\qty{286}{\celsius}の飽和蒸気をタービンに送ります。原子炉の中では燃料をいつも水でひたして冷やすので、出てくる蒸気は過熱されていません。過熱しないので $T_{\mathrm m}$ が低く、\qty{5}{kPa}の復水器でも理想の効率は約38\,\%、実際は約33\,\% です。同じ電気を出すのに、火力より多くの熱を海へ捨てることになります。また、タービンの出口の乾き度が0.67と小さいので、タービンのとちゅうで水の粒を取りのぞく装置（湿分分離器）をつけています。

\begin{phcode}[Python]
from iapws import IAPWS97          \# IAPWS-IF97 の式（P は MPa、T は K）
def rankine(pb, tb, pc=0.01):      \# ボイラーの圧力・温度、復水器の圧力
    s1 = IAPWS97(P=pb, T=tb+273.15)          \# 1：タービンの入口
    s2 = IAPWS97(P=pc, s=s1.s)               \# 2：s が同じ出口
    s3 = IAPWS97(P=pc, x=0)                  \# 3：飽和液
    w\_in = s3.v * (pb - pc) * 1000          \# ポンプ（kJ/kg）
    q\_in = s1.h - (s3.h + w\_in)
    w\_out = s1.h - s2.h
    return (w\_out - w\_in) / q\_in, s2.x
for tb in (500, 600):
    for pb in (2, 4, 10, 20):
        eta, x2 = rankine(pb, tb)
        print(tb, pb, round(eta, 3), round(x2, 3))
\end{phcode}
\begin{phcodeout}
500 2 0.339 0.905
500 4 0.368 0.859
500 10 0.402 0.793
500 20 0.421 0.733
600 2 0.357 0.941
600 4 0.384 0.896
600 10 0.417 0.834
600 20 0.438 0.781
\end{phcodeout}
\noindent ボイラーの圧力を上げると効率は上がりますが、出口の乾き度は下がります。温度を\qty{500}{\celsius}から\qty{600}{\celsius}に上げると、効率も乾き度も上がります。だから発電所は、高い圧力と高い温度と再熱を組み合わせています。
\begin{bkfigure}
\begin{tikzpicture}
\begin{axis}[bkplot, width=100mm, height=50mm, xmin=0, xmax=22, ymin=0.32, ymax=0.46, xtick={2,4,...,20}, ytick={0.34,0.38,0.42,0.46}, yticklabels={34,38,42,46}, xlabel={圧力（MPa）}, ylabel={効率（\%）}]
  \addplot[bkcurve, mark=*, mark size=1.4pt] coordinates {(2,0.339) (4,0.368) (6,0.384) (8,0.394) (10,0.402) (15,0.414) (20,0.421)};
  \addplot[bkcurve2, mark=square*, mark size=1.4pt] coordinates {(2,0.357) (4,0.384) (6,0.399) (8,0.410) (10,0.417) (15,0.430) (20,0.438)};
  \node[font=\sffamily\footnotesize, anchor=north west] at (axis cs:12,0.404) {\qty{500}{\celsius}（丸）};
  \node[font=\sffamily\footnotesize, anchor=south east] at (axis cs:11,0.43) {\qty{600}{\celsius}（四角）};
\end{axis}
\end{tikzpicture}
\bkfigcaption{理想のランキンサイクルの効率（縦）と、ボイラーの圧力（横。復水器\qty{10}{kPa}。プログラムで6・8・15 MPa も計算した）。}
\label{ph:熱-8:fig-sim}
\end{bkfigure}


\begin{phmada}[蒸気をもっと熱くできないのはなぜ？]
⑦のように、蒸気の温度を上げるほど、ランキンサイクルの効率は上がります。それなのに、最新の発電所の蒸気が\qty{600}{\celsius}くらいで止まっているのは、その温度と\qty{25}{MPa}をこえる圧力に、何十年も耐える管の材料が限られるからです。今の発電所のボイラーやタービンに使う鉄の合金（高クロム鋼）は、約\qty{630}{\celsius}が限界とされています。

日本では2008年から、蒸気を\qty{700}{\celsius}級にして、送り出す電気の効率を46\,\%（高位発熱量で数えた値）にする「先進超々臨界圧（A-USC）」の開発が進められ、ニッケルを多くふくむ合金で、ボイラーやタービンの部品の試験が行われました。ただし、ニッケルの合金は高価で、大きな部品を均一につくったり、溶接したりするのが難しく、何十年も使ったときに材料の中がどう変わるかも、試験の時間より長い分は予想にたよっています。こうした材料と費用の課題があり、実際の発電所で使われるには至っていません。
\end{phmada}

\begin{motto}[4]
\phboxhead{水のかわりに、わきやすい液体で回す}
水は、蒸発熱が大きく、安く、燃えず、毒もないので、ランキンサイクルに向いた液体です。ところが、温泉の湯や工場の排熱のように\qtyrange{80}{150}{\celsius}くらいの熱しかないと、水をわかしても圧力が低く、体積がとても大きな蒸気になって、タービンが大きくなりすぎます。そこで、水よりずっと低い温度でわく液体（ペンタンなどの炭化水素や、冷媒の仲間）を使うことがあります。これを\textbf{有機ランキンサイクル}（バイナリー発電）といい、熱い湯で別の液体をわかして、その蒸気でタービンを回します。温度差が小さいので、効率は1割くらいと低いのですが、捨てるはずだった熱から電気をつくれます。どの液体を選ぶかは、熱-9 の冷媒の選び方（燃えやすさ・地球温暖化への影響・圧力）とよく似た考え方で決めます。
\end{motto}

\begin{column}{ワットの復水器から、蒸気タービンへ}
蒸気の力を使う機械は、18世紀のイギリスで、炭鉱にたまる水をくみ上げるポンプとして広まりました（ニューコメンの機関、1712年）。この機関は、シリンダーの中に水をふきこんで蒸気を冷やし、ピストンを動かしていたので、1回ごとにシリンダーを冷やしては温め直し、多くの燃料をむだにしていました。ワットは、蒸気をシリンダーとは別の容器で冷やせば、シリンダーは熱いままでよいことに気づき、1769年にこの「分離凝縮器」の特許をとりました。これが、この章の復水器の始まりです。

19世紀には、ピストンを往復させる蒸気機関が、工場・機関車・船を動かしました。1884年、イギリスのパーソンズが、蒸気で羽根車を直接回す蒸気タービンを発明すると、速く回り、大きな出力を出せることから、発電所の主役は蒸気タービンになりました。ランキンが1859年の本で示した計算のしかたは、こうした機械の効率をくらべ、よりよい設計を選ぶための道具になりました。今の発電所の蒸気タービンも、この章と同じ蒸気表とバランス式で設計されています。
\end{column}

\begin{taisetsu}
\begin{itemize}[leftmargin=1.2\zw, itemsep=2pt, topsep=0pt, label={\color{bkNavy}●}]
\item 水は、飽和蒸気圧 $p_{\mathrm s}(T)$ が周りの圧力と等しくなる温度でわく。1気圧で\qty{99.97}{\celsius}、富士山頂で約\qty{87}{\celsius}。
\item 湿り蒸気：$h=h'+x(h''-h')$（$v$・$s$ も同じ形）。蒸発熱 $L=h''-h'$。表のあいだは直線補間。
\item ランキンサイクル：ボイラー $q_{\text{入}}=h_1-h_4$、タービン $w_{\text{出}}=h_1-h_2$、復水器 $q_{\text{出}}=h_2-h_3$、ポンプ $w_{\text{入}}\approx v_3(p_4-p_3)$。効率 $\eta=1-q_{\text{出}}/q_{\text{入}}$。
\item 効率を上げる：熱を受けとる平均の温度を上げる（高圧・高温・再熱・再生）か、捨てる温度を下げる。
\end{itemize}
\end{taisetsu}

\begin{kotoba}\small
\bkword{飽和蒸気圧}{液体とその蒸気がつり合ってたまるときの蒸気の圧力。温度だけで決まる。}
\bkword{沸騰}{飽和蒸気圧が周りの圧力に届き、液の中から泡ができてさかんに蒸発すること。}
\bkword{三重点・臨界点}{固体・液体・気体がいっしょにいられる点。液体と気体の区別がなくなる点。}
\bkword{飽和液・飽和蒸気}{ちょうどわく温度の液体（$'$ の量）、全部が蒸気になったばかりの蒸気（$''$ の量）。}
\bkword{湿り蒸気・過熱蒸気}{液と蒸気が混ざった状態。飽和温度より温度が高い蒸気。}
\bkword{乾き度}{湿り蒸気のうち、蒸気の質量の割合 $x$。}
\bkword{ランキンサイクル}{ボイラー・タービン・復水器・ポンプで水を回す、蒸気の熱機関のサイクル。}
\bkword{再熱・再生}{とちゅうの蒸気を温め直すこと。抽気した蒸気で給水を温めること。}
\end{kotoba}
\bkkigocards{\bkkigocard{$x$}{エックス}{乾き度。湿り蒸気のうち蒸気の質量の割合（0〜1）。熱-9 の絶対湿度 $x$ とは別の量（章が違う）}{dryness fraction, quality}{熱-8}}{}

\begin{senmonsho}[Rankine cycle]
\begin{bkhonbun}
For a saturated mixture, the enthalpy is $h=h_f+x\,h_{fg}$, where $h_{fg}=h_g-h_f$ is the enthalpy of vaporization. The thermal efficiency of the ideal Rankine cycle is
\[ \eta_{\mathrm{th}}=\frac{w_{\mathrm{net}}}{q_{\mathrm{in}}}=1-\frac{q_{\mathrm{out}}}{q_{\mathrm{in}}},\qquad w_{\mathrm{pump,in}}=v_1(P_2-P_1). \]
\end{bkhonbun}
\bkyaku{英語の本の書き方。$f$ は fluid（液）、$g$ は gas（蒸気）で、この本の $h'$ が $h_f$、$h''$ が $h_g$。$h_{fg}$ は蒸発熱 $L$。quality は乾き度。$q_{\mathrm{in}}$ や $w_{\mathrm{pump,in}}$ のように、この本と同じく「入・出」を添字で書く本が多い。状態の番号のつけ方（どこを1にするか）は本ごとに違うので、装置の図で確かめてから読む。}
\end{senmonsho}

%% ---------------- 例題 ----------------
\bksection{reidai}{例題でたしかめよう}

\noindent この章の例題は、蒸気表（表\ref{ph:熱-8:tab-ps}〜表\ref{ph:熱-8:tab-kanetsu}、巻末の蒸気表）の値を使うので、とくにことわりがなければ答えを3けたで書きます。圧力はすべて絶対圧です。

\begin{phreidai}[富士山の頂上で、水は何℃でわく？]
\label{ph:熱-8:ex1}富士山の頂上の気圧は約\qty{63}{kPa}です。巻末の蒸気表には、\qty{85}{\celsius}で $p_{\mathrm s}=\qty{57.867}{kPa}$、\qty{90}{\celsius}で $p_{\mathrm s}=\qty{70.182}{kPa}$ とあります。水がわく温度を求めましょう（答えは2けた）。
\phbamen
\bkhatenaref の場面です。なべの水の上には、約\qty{63}{kPa}の空気が乗っています。
\phmushi
なべの水の深さの分の圧力（数 cm の水で\qty{0.5}{kPa}くらい）と、なべの中の温度のむらを無視します。気圧は天気でも少し変わりますが、\qty{63}{kPa}とします。
\phzu
周りの圧力 $p=\qty{63}{kPa}$、わく温度 $\theta$〔℃〕。表の2行を $(\theta_1,\ p_1)=(85,\ 57.867)$、$(\theta_2,\ p_2)=(90,\ 70.182)$。
\phriyu
水は $p_{\mathrm s}=p$ となる温度でわく（①）。表では $p_{\mathrm s}$ から温度を逆に読むので、温度を圧力で直線補間する（③）。5 K の間かくなら、曲がりは小さい。
\phtateru
\[ \theta=\theta_1+(\theta_2-\theta_1)\times\frac{p-p_1}{p_2-p_1} \]
\phtoku
\[ \frac{p-p_1}{p_2-p_1}=\frac{63-57.867}{70.182-57.867}=\frac{5.133}{12.315}\approx0.417 \]
\[ \theta=85+5\times0.417\approx\qty{87.1}{\celsius} \]
\phtashikame
単位：℃＋K×（無次元）＝℃。大きさ：$p$ は\qty{57.9}{kPa}と\qty{70.2}{kPa}の間なので、$\theta$ も\qty{85}{\celsius}と\qty{90}{\celsius}の間に入る。IF97 の式で詳しく解くと\qty{87.18}{\celsius}で、補間との差は\qty{0.1}{K}くらい。\qty{80}{\celsius}と\qty{100}{\celsius}の行で補間すると\qty{85.8}{\celsius}と、\qty{1.4}{K}もずれる（曲がった線を、長い直線でつないだため）。極端な場合：$p=\qty{101.325}{kPa}$ なら\qty{99.97}{\celsius}にもどる。
\kotae 約\qty{87}{\celsius}
\end{phreidai}
\begin{jibun}
エベレストの頂上の気圧は約\qty{34}{kPa}です。巻末の蒸気表の\qty{70}{\celsius}（\qty{31.201}{kPa}）と\qty{75}{\celsius}（\qty{38.595}{kPa}）の行から、水がわく温度を求めましょう（2けた）。
\bkjans{$34-31.201=2.799$、$38.595-31.201=7.394$、$2.799\div7.394\approx0.379$。$70+5\times0.379\approx71.9$ で、約\qty{72}{\celsius}}
\end{jibun}

\begin{reidai}[蒸気表を読む]
\label{ph:熱-8:ex2}\qty{1}{MPa}の飽和の状態について、表\ref{ph:熱-8:tab-sat}から、わく温度、飽和液と飽和蒸気の比エンタルピー、蒸発熱を読みましょう。
\par\bkstep{理想化} 表の値は、純粋な水の、つり合った状態の値（不純物や流れは考えない）。
\par\kangaekata 圧力の列が1の行を、横に読む。蒸発熱は $L=h''-h'$（②）。
\par\shiki $\theta=\qty{179.89}{\celsius}$、$h'=\qty{762.7}{kJ/kg}$、$h''=\qty{2777.1}{kJ/kg}$、$L=2777.1-762.7$
\par\bkstep{確かめ} 1気圧の $L$（\qty{2256.5}{kJ/kg}）より小さい。圧力（わく温度）が高いほど $L$ は小さくなる（表\ref{ph:熱-8:tab-ps}）ので、合っている。
\par\kotae \qty{179.89}{\celsius}、$h'=\qty{762.7}{kJ/kg}$、$h''=\qty{2777.1}{kJ/kg}$、$L=\qty{2014.4}{kJ/kg}$
\end{reidai}
\begin{jibun}
\qty{0.5}{MPa}の飽和の状態について、表\ref{ph:熱-8:tab-sat}から、わく温度、$h'$、$h''$、蒸発熱を読みましょう。
\bkjans{\qty{151.84}{\celsius}、$h'=\qty{640.2}{kJ/kg}$、$h''=\qty{2748.1}{kJ/kg}$、$L=\qty{2107.9}{kJ/kg}$}
\end{jibun}

\begin{phreidai}[表のあいだを直線で読む：直線補間のずれ]
\label{ph:熱-8:ex3}表\ref{ph:熱-8:tab-sat}の\qty{100}{kPa}と\qty{200}{kPa}の行から、\qty{150}{kPa}の飽和温度を直線補間で求めましょう。本当の値（巻末の表）は\qty{111.35}{\celsius}です。どれだけずれるか、ずれるわけも考えましょう。
\phbamen
表にない圧力の飽和温度を、表の2行から読む場面です。
\phmushi
2行の間の飽和温度の線を、まっすぐとみなします（これが、ずれの原因になる）。
\phzu
$(p_1,\ \theta_1)=(100,\ 99.61)$、$(p_2,\ \theta_2)=(200,\ 120.21)$、$p=\qty{150}{kPa}$（温度は℃）。
\phriyu
3-5⑧の直線内挿。$p$ がちょうど2行のまん中なので、割合は0.5。
\phtateru
\[ \theta=\theta_1+(\theta_2-\theta_1)\times\frac{p-p_1}{p_2-p_1} \]
\phtoku
\[ \theta=99.61+(120.21-99.61)\times\frac{150-100}{200-100}=99.61+20.60\times0.5\approx\qty{109.91}{\celsius} \]
本当の値との差は $111.35-109.91=\qty{1.44}{K}$ です。
\phtashikame
大きさ：2つの温度の間に入っている。ずれの向き：飽和温度は、圧力に対して上に凸の曲線（圧力が倍になっても、温度は倍にならず、増え方がゆるやかになる。図\ref{ph:熱-8:fig-sozu}）なので、2点を結ぶ直線は、いつも曲線の下を通る。だから補間の値は本当の値より低くなる。表の間かくを\qty{10}{kPa}くらいに細かくすれば、ずれは0.1 K 以下になる。
\kotae 約\qty{109.9}{\celsius}（本当の値より約\qty{1.4}{K}低い）
\end{phreidai}
\begin{jibun}
表\ref{ph:熱-8:tab-sat}の\qty{1}{MPa}と\qty{2}{MPa}の行から、\qty{1.5}{MPa}の飽和温度を直線補間で求めましょう。本当の値（\qty{198.30}{\celsius}）とくらべて、どちらにずれますか。
\bkjans{$212.38-179.89=32.49$、$179.89+32.49\times0.5\approx\qty{196.1}{\celsius}$。本当の値より約\qty{2.2}{K}低い（上に凸の曲線の下を直線が通る）}
\end{jibun}

\begin{reidai}[湿り蒸気の比エンタルピー]
\label{ph:熱-8:ex4}\qty{10}{kPa}・乾き度0.90の湿り蒸気の比エンタルピーを求めましょう。
\par\bkstep{理想化} 水と水蒸気はどちらも\qty{10}{kPa}の飽和の状態にあり、つり合っているとする。
\par\kangaekata $h=h'+x(h''-h')$（②）。表\ref{ph:熱-8:tab-sat}の\qty{10}{kPa}（\qty{0.01}{MPa}）の行で、$h'=\qty{191.8}{kJ/kg}$、$h''=\qty{2583.9}{kJ/kg}$。
\par\shiki $h=191.8+0.90\times(2583.9-191.8)$\\ \hspace*{2\zw}$=191.8+0.90\times2392.1=191.8+2152.9$
\par\bkstep{確かめ} $h'$ と $h''$ の間で、$h''$ に近い（9割が蒸気）。
\par\kotae 約\qty{2345}{kJ/kg}（\qty{2344.7}{kJ/kg}）
\end{reidai}
\begin{jibun}
\qty{100}{kPa}・乾き度0.95の湿り蒸気の比エンタルピーを、表\ref{ph:熱-8:tab-sat}から求めましょう。
\bkjans{$2674.9-417.4=2257.5$、$417.4+0.95\times2257.5$ $\approx\qty{2562}{kJ/kg}$}
\end{jibun}

\begin{reidai}[やかんの水は、1分で何 g 蒸発する？]
\label{ph:熱-8:ex5}1気圧で、わいているやかんに、ヒーターが1秒あたり\qty{1.0}{kJ}（\qty{1.0}{kW}）の熱を加え続けます。1分間に蒸発する水の質量を、表\ref{ph:熱-8:tab-ps}の\qty{100}{\celsius}の蒸発熱から求めましょう。
\par\bkstep{理想化} ヒーターの熱は全部水に入り、やかんから周りへ逃げる熱は無視する。水はずっと\qty{100}{\celsius}でわいているとする（\qty{99.97}{\celsius}との差は無視）。
\par\kangaekata わいている間、加えた熱は全部蒸発に使われる。蒸発する質量＝熱÷蒸発熱（$Q=mL$、熱-1）。
\par\shiki $Q=\qty{1.0}{kW}\times\qty{60}{s}=\qty{60}{kJ}$、$m=60\div2256.5\approx\qty{0.0266}{kg}$
\par\bkstep{確かめ} 単位：kJ÷（kJ/kg）＝kg。大きさ：1分で大さじ2杯くらい。やかんの水\qty{1}{L}がなくなるまでには約38分かかり、空だきに気づく時間がある。
\par\kotae 約\qty{26.6}{g}
\end{reidai}
\begin{jibun}
\qty{200}{\celsius}でわいている水（ボイラーの中）に、\qty{60}{kJ}の熱を加えると、何 g 蒸発しますか。表\ref{ph:熱-8:tab-ps}の $L$ を使いましょう。
\bkjans{$60\div1939.7\approx\qty{0.0309}{kg}$。約\qty{30.9}{g}（\qty{100}{\celsius}のときより多い。高い温度ほど $L$ が小さい）}
\end{jibun}

\begin{phreidai}[ランキンサイクルの効率を計算する]
\label{ph:熱-8:ex6}ボイラーで\qty{10}{MPa}・\qty{500}{\celsius}の蒸気をつくり、タービンで\qty{10}{kPa}まで膨張させ、復水器で飽和液にもどす理想のランキンサイクルがあります。タービンの仕事・ポンプの仕事・ボイラーで受けとる熱と、このサイクルの効率を求めましょう。
\phbamen
火力発電所の基本の形です（図\ref{ph:熱-8:fig-sochi}・図\ref{ph:熱-8:fig-ts}）。
\phmushi
⑥の「理想のランキンサイクル」のモデル（管の摩擦、周りへの熱、速さ・高さの差を無視。タービン・ポンプは断熱・可逆）を使います。
\phzu
状態1（タービンの入口：\qty{10}{MPa}・\qty{500}{\celsius}）、2（タービンの出口：\qty{10}{kPa}・$s_2=s_1$）、3（復水器の出口：\qty{10}{kPa}の飽和液）、4（ポンプの出口：\qty{10}{MPa}）。
\phriyu
4つの機械に熱-4 のバランス式。タービンの出口は $s$ が一定なので、$s$ から乾き度を出す。ポンプは水が縮まないので $v\Delta p$。
\phtateru
表\ref{ph:熱-8:tab-kanetsu}から $h_1=\qty{3375.1}{kJ/kg}$、$s_1=\qty{6.5993}{kJ/(kg.K)}$。表\ref{ph:熱-8:tab-sat}の\qty{10}{kPa}の行から $h'=191.8$、$h''=2583.9$、$s'=0.6492$、$s''=8.1489$、$v'=0.001010$。
\[ x_2=\frac{s_1-s'}{s''-s'},\qquad h_2=h'+x_2(h''-h') \]
\[ w_{\text{入}}=v_3(p_4-p_3),\qquad \eta=\frac{w_{\text{出}}-w_{\text{入}}}{q_{\text{入}}} \]
\phtoku
タービンの出口：
\[ x_2=\frac{6.5993-0.6492}{8.1489-0.6492}=\frac{5.9501}{7.4997}\approx0.7934 \]
\[ h_2=191.8+0.7934\times2392.1\approx\qty{2089.6}{kJ/kg} \]
タービン：$w_{\text{出}}=3375.1-2089.6=\qty{1285.5}{kJ/kg}$。

ポンプ（kPa×m$^3$/kg＝kJ/kg）：
\[ w_{\text{入}}=0.001010\times(\num{10000}-10)\approx\qty{10.1}{kJ/kg} \]
$h_4=191.8+10.1=\qty{201.9}{kJ/kg}$。

ボイラー：$q_{\text{入}}=3375.1-201.9=\qty{3173.2}{kJ/kg}$。

効率：
\[ \eta=\frac{1285.5-10.1}{3173.2}=\frac{1275.4}{3173.2}\approx0.402 \]
\phtashikame
単位：kJ/kg どうしの比なので、単位なし。\\
バランス：$q_{\text{出}}=h_2-h_3=2089.6-191.8=\qty{1897.8}{kJ/kg}$ で、
\[ q_{\text{入}}+w_{\text{入}}=3183.3,\qquad w_{\text{出}}+q_{\text{出}}=3183.3 \]
と等しい。大きさ：カルノーの効率（\qty{773}{K}と\qty{319}{K}で0.59）より小さく、平均の受熱温度で考えた0.40（⑥）と同じ。極端な場合：復水器の圧力を上げてボイラーと同じにすると、タービンは仕事をせず $\eta=0$。
\kotae タービン\qty{1286}{kJ/kg}、ポンプ\qty{10.1}{kJ/kg}、ボイラー\qty{3173}{kJ/kg}、効率 約40.2\,\%
\end{phreidai}
\begin{jibun}
\bkreidairef{6}のサイクルで、タービンの断熱効率が0.88だとすると（⑧）、タービンの仕事、出口の乾き度、サイクルの効率はいくらになりますか。ポンプとボイラーは\bkreidairef{6}と同じとします。
\bkjans{$w_{\text{出}}=0.88\times1285.5$ $\approx\qty{1131.2}{kJ/kg}$。$h_2=\qty{2243.9}{kJ/kg}$。$x_2=2052.1\div2392.1$ $\approx0.858$。$\eta=1121.1\div3173.2$ $\approx0.353$（約35.3\,\%）}
\end{jibun}

\begin{reidai}[ポンプの仕事はどれくらい？]
\label{ph:熱-8:ex7}\bkreidairef{6}のポンプの仕事は、タービンの仕事の何\,\% ですか。
\par\bkstep{理想化} 水は縮まず、ポンプは断熱・可逆とする（\bkreidairef{6}と同じ）。
\par\kangaekata 2つの仕事の比をとる。$w_{\text{入}}=v_3(p_4-p_3)$。
\par\shiki $10.1\div1285.5\approx0.0079$
\par\bkstep{確かめ} 水の比体積（約\qty{0.001}{m^3/kg}）は、タービンの中の蒸気（\qty{0.03}{m^3/kg}〜\qty{12}{m^3/kg}）よりずっと小さいので、同じ圧力の差でも仕事が小さい。熱-4 の空気の圧縮機（1 kg あたり数百 kJ）とくらべても、とても小さい。
\par\kotae 約0.79\,\%（1\,\% にもならない）
\end{reidai}
\begin{jibun}
\qty{10}{kPa}の飽和液（$v'=\qty{0.001010}{m^3/kg}$）を、ポンプで\qty{15}{MPa}まで押すときの1 kg あたりの仕事を求めましょう。
\bkjans{$0.001010\times(\num{15000}-10)$ $\approx\qty{15.1}{kJ/kg}$}
\end{jibun}

\begin{reidai}[タービンの出口の蒸気は、どれだけふくらむ？]
\label{ph:熱-8:ex8}\bkreidairef{6}のタービンの出口（\qty{10}{kPa}・$x_2=0.7934$）の比体積を求め、入口（表\ref{ph:熱-8:tab-kanetsu}で $v_1=\qty{0.03281}{m^3/kg}$）の何倍かを求めましょう。
\par\bkstep{理想化} 出口の水の粒と蒸気は、つり合った飽和の状態とする。
\par\kangaekata $v=v'+x(v''-v')$（②）。\qty{10}{kPa}の $v'=0.001010$、$v''=14.67$。
\par\shiki $v_2=0.001010+0.7934\times(14.67-0.001010)$\\ \hspace*{2\zw}$\approx\qty{11.64}{m^3/kg}$、$11.64\div0.03281\approx355$
\par\bkstep{確かめ} 出口は2割が水なので、飽和蒸気の $v''$ の約8割。1秒に流れる質量は同じなので、出口では体積が約355倍になる。低圧のタービンの羽根が、とても長い（\qty{1}{m}をこえる）理由である。
\par\kotae 約\qty{11.6}{m^3/kg}、入口の約355倍
\end{reidai}
\begin{jibun}
\qty{4}{MPa}・\qty{500}{\celsius}の蒸気（表\ref{ph:熱-8:tab-kanetsu}）を、\qty{10}{kPa}まで断熱・可逆に膨張させます。出口の乾き度と比体積を求め、入口の比体積の何倍か求めましょう。
\bkjans{$x=6.4427\div7.4997\approx0.859$。$v\approx0.001010+0.859\times14.669$ $\approx\qty{12.6}{m^3/kg}$。$12.6\div0.08644\approx146$ 倍}
\end{jibun}

\begin{phreidai}[再熱サイクルの効率]
\label{ph:熱-8:ex9}\bkreidairef{6}のサイクルで、高圧のタービンで\qty{4}{MPa}まで膨張した蒸気を、ボイラーにもどして\qty{500}{\celsius}まで温め直し（再熱）、低圧のタービンで\qty{10}{kPa}まで膨張させます。効率と、低圧のタービンの出口の乾き度を求めましょう。\qty{4}{MPa}まで膨張した蒸気は、IF97 の式で計算すると、$s=\qty{6.5993}{kJ/(kg.K)}$ で\qty{353.8}{\celsius}・$h_{\mathrm a}=\qty{3102.7}{kJ/kg}$ です。
\phbamen
タービンの出口の水の粒を減らすくふうです（図\ref{ph:熱-8:fig-saisei}の「再熱」）。
\phmushi
\bkreidairef{6}と同じ理想化。再熱器の中でも圧力は\qty{4}{MPa}で一定とします。
\phzu
1（\qty{10}{MPa}・\qty{500}{\celsius}）→ a（\qty{4}{MPa}、$s=s_1$）→ b（\qty{4}{MPa}・\qty{500}{\celsius}）→ c（\qty{10}{kPa}、$s=s_{\mathrm b}$）→ 3 → 4。
\phriyu
タービンの仕事は2つの部分の和、ボイラーの熱は「4→1」と「a→b」の和。c の乾き度は $s$ から出す。
\phtateru
\[ w_{\text{出}}=(h_1-h_{\mathrm a})+(h_{\mathrm b}-h_{\mathrm c}),\qquad q_{\text{入}}=(h_1-h_4)+(h_{\mathrm b}-h_{\mathrm a}) \]
\phtoku
表\ref{ph:熱-8:tab-kanetsu}から $h_{\mathrm b}=\qty{3445.8}{kJ/kg}$、$s_{\mathrm b}=\qty{7.0919}{kJ/(kg.K)}$。
\[ x_{\mathrm c}=\frac{7.0919-0.6492}{7.4997}\approx0.8591 \]
\[ h_{\mathrm c}=191.8+0.8591\times2392.1\approx\qty{2246.8}{kJ/kg} \]
\[ w_{\text{出}}=(3375.1-3102.7)+(3445.8-2246.8)=272.4+1199.0=\qty{1471.4}{kJ/kg} \]
\[ q_{\text{入}}=(3375.1-201.9)+(3445.8-3102.7)=3173.2+343.1=\qty{3516.3}{kJ/kg} \]
\[ \eta=\frac{1471.4-10.1}{3516.3}\approx0.416 \]
\phtashikame
単位：比なので単位なし。大きさ：再熱なし（0.402・$x=0.793$）より、効率は少し上がり、乾き度は大きく上がった。再熱の熱343 kJ/kg は、\qty{354}{\celsius}から\qty{500}{\celsius}の高い温度で受けとる熱なので、平均の受熱温度が上がる。極端な場合：再熱の圧力を\qty{10}{MPa}（入口と同じ）にすると、再熱しないのと同じになる。
\kotae 効率 約41.6\,\%、出口の乾き度 約0.859
\end{phreidai}
\begin{jibun}
\bkreidairef{9}で、再熱の圧力を\qty{2}{MPa}にすると、どうなりますか。\qty{2}{MPa}まで膨張した蒸気は $h_{\mathrm a}=\qty{2930.7}{kJ/kg}$、\qty{2}{MPa}・\qty{500}{\celsius}の蒸気は $h_{\mathrm b}=\qty{3468.1}{kJ/kg}$・$s_{\mathrm b}=\qty{7.4335}{kJ/(kg.K)}$ です。
\bkjans{$x_{\mathrm c}=6.7843\div7.4997\approx0.905$、$h_{\mathrm c}\approx\qty{2355.7}{kJ/kg}$。$w_{\text{出}}=444.4+1112.4$ $=\qty{1556.8}{kJ/kg}$。$q_{\text{入}}=3173.2+537.4$ $=\qty{3710.6}{kJ/kg}$。$\eta=1546.7\div3710.6$ $\approx0.417$}
\end{jibun}

\begin{phreidai}[再生サイクル：抽気で給水を温める]
\label{ph:熱-8:ex10}\bkreidairef{6}のサイクルで、タービンのとちゅうの\qty{0.5}{MPa}から蒸気を取り出し、開放形の給水加熱器で給水と混ぜて、\qty{0.5}{MPa}の飽和液にしてからボイラーへ送ります。抽気の割合と効率を求めましょう。
\phbamen
図\ref{ph:熱-8:fig-saisei}の「再生」の装置です。ポンプは2台（\qty{10}{kPa}→\qty{0.5}{MPa}と、\qty{0.5}{MPa}→\qty{10}{MPa}）になります。
\phmushi
\bkreidairef{6}と同じ理想化。加熱器から周りへ逃げる熱は無視します。
\phzu
タービンに入る蒸気を1 kg とし、抽気を $y$〔kg〕とします。6：抽気（\qty{0.5}{MPa}、$s=s_1$）、3$'$：1台めのポンプの出口、f：\qty{0.5}{MPa}の飽和液、4：2台めのポンプの出口。
\phriyu
加熱器は「入るエネルギー＝出るエネルギー」の混合（熱-4）。タービンの仕事は、1 kg が6まで、$(1-y)$〔kg〕が6から2まで膨張する分。
\phtateru
\[ y\,h_6+(1-y)\,h_{3'}=h_{\mathrm f}\quad\text{より}\quad y=\frac{h_{\mathrm f}-h_{3'}}{h_6-h_{3'}} \]
\[ w_{\text{出}}=(h_1-h_6)+(1-y)(h_6-h_2),\qquad q_{\text{入}}=h_1-h_4 \]
\phtoku
6：\qty{0.5}{MPa}の $s''=6.8206$ より $s_1$ が小さいので湿り蒸気。
\[ x_6=\frac{6.5993-1.8606}{6.8206-1.8606}\approx0.9554 \]
\[ h_6=640.2+0.9554\times2107.9\approx\qty{2654.1}{kJ/kg} \]
ポンプは、1台めと2台めで
\[ 0.001010\times490\approx\qty{0.49}{kJ/kg},\qquad 0.001093\times9500\approx\qty{10.4}{kJ/kg} \]
なので、$h_{3'}=\qty{192.3}{kJ/kg}$、$h_4=\qty{650.6}{kJ/kg}$。
\[ y=\frac{640.2-192.3}{2654.1-192.3}=\frac{447.9}{2461.8}\approx0.182 \]
\[ w_{\text{出}}=(3375.1-2654.1)+0.818\times(2654.1-2089.6) \]
\[ \approx721.0+461.8=\qty{1182.8}{kJ/kg} \]
ポンプの仕事の合計とボイラーの熱は
\[ 0.818\times0.49+10.4\approx\qty{10.8}{kJ/kg},\qquad q_{\text{入}}=3375.1-650.6=\qty{2724.5}{kJ/kg} \]
です。
\[ \eta=\frac{1182.8-10.8}{2724.5}\approx0.430 \]
\phtashikame
大きさ：タービンの仕事は再生なしより減った（\qty{1285.5}{kJ/kg}→\qty{1182.8}{kJ/kg}）が、ボイラーの熱がもっと減った（給水が\qty{152}{\celsius}まで温まっている）ので、効率は0.402から0.430に上がった。極端な場合：$y=0$（抽気しない）なら\bkreidairef{6}にもどる。
\kotae 抽気の割合 約0.18、効率 約43.0\,\%
\end{phreidai}
\begin{jibun}
\bkreidairef{10}で、抽気の圧力を\qty{1}{MPa}にします。抽気の割合と効率を求めましょう。使う値は次のとおりです。
\par\noindent ・\qty{1}{MPa}で $s=s_1$ の蒸気：$h_6=\qty{2783.6}{kJ/kg}$
\par\noindent ・\qty{1}{MPa}の飽和液：$h'=\qty{762.7}{kJ/kg}$、$v'=\qty{0.001127}{m^3/kg}$
\par\noindent ・1台めのポンプの出口：$h_{3'}=\qty{192.8}{kJ/kg}$
\bkjans{$y=569.9\div2590.8\approx0.220$。$w_{\text{出}}=591.5+0.780\times694.0$ $\approx\qty{1132.8}{kJ/kg}$。ポンプ $0.78+10.14\approx\qty{10.9}{kJ/kg}$。$q_{\text{入}}=3375.1-772.8$ $=\qty{2602.3}{kJ/kg}$。$\eta\approx0.431$}
\end{jibun}

\begin{reidai}[100万 kW の発電所が捨てる熱]
\label{ph:熱-8:ex11}\bkreidairef{6}のサイクルで、電気を100万 kW（\qty{1000}{MW}）出すには、1秒に何 kg の蒸気を回せばよいですか。復水器で捨てる熱は何 MW ですか。
\par\bkstep{理想化} 正味の仕事が全部電気になる（発電機などの損失は無視）とする。
\par\kangaekata 1秒に $\dot m$〔kg〕回すと、出力は $\dot m\,(w_{\text{出}}-w_{\text{入}})$、捨てる熱は $\dot m\,q_{\text{出}}$。
\par\shiki $w_{\text{出}}-w_{\text{入}}=1285.5-10.1=\qty{1275.4}{kJ/kg}$。\\ \hspace*{2\zw}$\dot m=\num{1000000}\div1275.4\approx\qty{784}{kg/s}$。捨てる熱 $784\times1897.8\approx\num{1488000}$\,kW
\par\bkstep{確かめ} 単位：kW÷（kJ/kg）＝kg/s。大きさ：捨てる熱は電気の約1.5倍で、熱-0 の「燃料の熱の半分以上が海へ」と合う。この熱を、温度差\qty{7}{K}の海水（\qty{1025}{kg/m^3}、\qty{3.99}{kJ/(kg.K)}）で運ぶと、1秒に約\qty{52}{m^3}の海水が要る（熱-0 の計算と同じ程度）。
\par\kotae 約\qty{784}{kg/s}、約\qty{1490}{MW}
\end{reidai}
\begin{jibun}
\bkreidairef{6}のサイクルで、電気を50万 kW 出す発電所の蒸気の量と、復水器で捨てる熱を求めましょう。
\bkjans{$\dot m=\num{500000}\div1275.4$ $\approx\qty{392}{kg/s}$。捨てる熱 $392\times1897.8$ $\approx\qty{744}{MW}$}
\end{jibun}

\begin{reidai}[過熱蒸気表を補間して読む]
\label{ph:熱-8:ex12}\qty{1}{MPa}・\qty{280}{\celsius}の過熱蒸気の比エンタルピーを、表\ref{ph:熱-8:tab-kanetsu}から直線補間で求めましょう。
\par\bkstep{理想化} \qty{250}{\celsius}と\qty{300}{\celsius}の間で、$h$ は温度にほぼ比例して増えるとする。
\par\kangaekata \qty{1}{MPa}の行の\qty{250}{\celsius}（\qty{2943.2}{kJ/kg}）と\qty{300}{\celsius}（\qty{3051.7}{kJ/kg}）を補間する。割合は $(280-250)\div(300-250)=0.6$。
\par\shiki $h=2943.2+(3051.7-2943.2)\times0.6=2943.2+65.1$
\par\bkstep{確かめ} IF97 の式で計算した値は\qty{3008.7}{kJ/kg}で、差は\qty{0.4}{kJ/kg}。過熱蒸気の $h$ は温度にほぼ比例するので、補間がよく当たる（飽和蒸気圧のような急な曲線ではない）。
\par\kotae 約\qty{3008}{kJ/kg}（\qty{3008.3}{kJ/kg}）
\end{reidai}
\begin{jibun}
\qty{4}{MPa}・\qty{450}{\celsius}の過熱蒸気の比エンタルピーを、表\ref{ph:熱-8:tab-kanetsu}から直線補間で求めましょう。
\bkjans{$(3214.4+3445.8)\div2$ $=\qty{3330.1}{kJ/kg}$（巻末の蒸気表では\qty{3331.0}{kJ/kg}）}
\end{jibun}

\begin{reidai}[圧力なべの中の温度]
\label{ph:熱-8:ex13}圧力なべの圧力計（ゲージ圧）が\qty{99}{kPa}を示しています。大気圧を\qty{101.3}{kPa}として、中の水は何℃でわいていますか。
\par\bkstep{理想化} なべの中は水と水蒸気だけで、空気は追い出されているとする。
\par\kangaekata 蒸気表は絶対圧で読む（力-12）。絶対圧＝ゲージ圧＋大気圧。
\par\shiki 絶対圧 $99+101.3=\qty{200.3}{kPa}$。表\ref{ph:熱-8:tab-sat}の\qty{0.2}{MPa}の行で\qty{120.21}{\celsius}
\par\bkstep{確かめ} ゲージ圧のまま読むと約\qty{99}{\celsius}で、ふつうのなべと変わらない。
\par\kotae 約\qty{120}{\celsius}
\end{reidai}
\begin{jibun}
ゲージ圧\qty{80}{kPa}の圧力なべ（大気圧\qty{101.3}{kPa}）の中で、水がわく温度を、表\ref{ph:熱-8:tab-sat}の\qty{0.1}{MPa}と\qty{0.2}{MPa}の行から直線補間で求めましょう（3けた）。
\bkjans{絶対圧\qty{181.3}{kPa}。割合 $81.3\div100=0.813$。$99.61+20.60\times0.813$ $\approx\qty{116}{\celsius}$（IF97 の式では\qty{117.1}{\celsius}。補間は少し低めに出る）}
\end{jibun}

%% ---------------- まちがえやすいところ ----------------
\bksection{tsumazuki}{まちがえやすいところをチェック}
\begin{tsumazuki}{ゲージ圧のまま蒸気表を読む}
\bkng{圧力計が\qty{1}{MPa}を示しているボイラーの蒸気は、表の\qty{1}{MPa}の行で読む。}
\bkok{ふつうの圧力計はゲージ圧。約\qty{0.1}{MPa}足した絶対圧\qty{1.1}{MPa}で読む。}
蒸気表は絶対圧で作られています。圧力が低いほど、この差が大きくききます（\bkreidairef{13}）。
\end{tsumazuki}
\begin{tsumazuki}{湯気は水蒸気？}
\bkng{やかんの口から出る白い湯気は、水蒸気である。}
\bkok{水蒸気は透明で、目に見えない。白い湯気は、水蒸気が冷えてできた小さな水の粒。}
やかんの口のすぐ先をよく見ると、何も見えないすき間があります。そこが水蒸気です。少しはなれて冷えると、水蒸気が水の粒にもどり、光を散らして白く見えます。空気の中で水蒸気が水にもどる条件は、熱-9 で調べます。
\end{tsumazuki}
\begin{tsumazuki}{湿り蒸気の $h$ を $xh''$ にする}
\bkng{乾き度0.9の蒸気の比エンタルピーは $0.9\times h''$。}
\bkok{$h=(1-x)h'+xh''=h'+x(h''-h')$。残りの1割の水の $h'$ も足す。}
$xh''$ だけでは、1 kg のうち水の\qty{0.1}{kg}の分のエネルギーが抜けてしまいます。$x=0$ のとき $h=h'$ になるか、で確かめましょう。
\end{tsumazuki}
\begin{tsumazuki}{効率の式に ℃ を入れる}
\bkng{\qty{500}{\celsius}と\qty{46}{\celsius}のカルノーの効率は $1-46/500\approx0.91$。}
\bkok{温度は K で入れる。$1-319/773\approx0.59$。}
温度の比をとる式（熱-5）は、絶対温度でしか成り立ちません。表の温度は℃なので、273.15 を足してから使います。
\end{tsumazuki}
\begin{tsumazuki}{ポンプの仕事を引き忘れる}
\bkng{ランキンサイクルの効率は $w_{\text{出}}/q_{\text{入}}$（タービンの仕事だけ）。}
\bkok{$\eta=(w_{\text{出}}-w_{\text{入}})/q_{\text{入}}$。ポンプに入れる仕事を引く。}
この章の $w_{\text{出}}$ はタービン1台の仕事です。熱-5 の $W_{\text{出}}$ や熱-7 の $w_{\text{出}}$（1周で外へ出す正味の仕事）にあたるのは、$w_{\text{出}}-w_{\text{入}}$ のほうです。ポンプの仕事は小さい（\bkreidairef{7}）ので、引き忘れても効率は少ししか変わりませんが、ガスタービン（熱-7）では圧縮機の仕事が大きいので、大きく間違えます。
\end{tsumazuki}

%% ---------------- 工学ではこう使う ----------------
\Needspace{16\baselineskip}
\bksection{phgear}{工学ではこう使う}
\begin{phkougaku}[発電所と工場の蒸気]
\phboxhead{① 発電所の蒸気の道すじ}
大きな火力発電所では、ボイラーで1秒に数百 kg の水を蒸気にし、高圧・中圧・低圧の3つのタービンに順に通します（とちゅうで再熱）。タービンの回る速さは、東日本（50 Hz）で毎分3000回転、西日本（60 Hz）で毎分3600回転と、電気の周波数で決まります。タービンの出口の蒸気は⑤のとおり体積がとても大きいので、低圧のタービンは何台も並べ、最後の羽根は長さ\qty{1}{m}をこえます。

\phboxhead{② 工場の蒸気の管とスチームトラップ}
工場で蒸気を管で送ると、とちゅうで熱が逃げて、蒸気の一部が水（ドレン）にもどります。たまった水は、管の中を蒸気の速さで流されて曲がり角に当たり、大きな音と力を出すことがあります（スチームハンマー）。そこで、管の低い所に、水だけを外へ出して蒸気は通さない弁（スチームトラップ）をつけます。管を保温材で包むのは、熱をむだにしないためだけでなく、水をつくらないためでもあります（熱-11）。

\phboxhead{③ 地熱発電とごみ発電}
地熱発電は、地下の熱い水を井戸でくみ上げ、圧力を下げてわかした蒸気（フラッシュ蒸気）でタービンを回します。温度がやや低いので、効率は火力より低くなりますが、燃料が要りません。ごみ焼却場の発電も、燃やした熱でボイラーの水をわかすランキンサイクルです。ごみにふくまれる塩素などで管がいたみやすいので、蒸気の温度を\qtyrange{400}{450}{\celsius}くらいにおさえることが多く、そのぶん効率は低めです。
\begin{phketa}[蒸気と発電所の数]
\phketaline{1気圧の水の蒸発熱}{約\qty{2.26}{MJ/kg}（\qty{1}{kg}を\qty{1}{K}温める熱の約540倍）}
\phketaline{\qty{100}{\celsius}で、わいてふくらむ体積}{約1600倍}
\phketaline{100万 kW の発電所の蒸気の量}{約\qtyrange{600}{800}{kg/s}}
\phketaline{復水器の圧力}{約\qty{5}{kPa}（大気圧の約20分の1）}
\phketaline{理想のランキンサイクルの効率}{\qtyrange{35}{48}{\percent}（条件による）}
\phboxhead{答えが桁で合っているか}
蒸気の $h$ は、ほとんどの場合\qtyrange{2000}{3700}{kJ/kg}の間、水の $h$ は\qty{1500}{kJ/kg}より小さい（臨界点の近く以外）。タービンの仕事が1 kg あたり\qty{2000}{kJ}をこえたり、ポンプの仕事が\qty{100}{kJ}をこえたりしたら、単位（kPa と MPa）や表の読み方を見直す。
\end{phketa}
\end{phkougaku}
\phstation{工業熱力学,蒸気動力・発電工学,ボイラー・蒸気タービン,原子力の熱水力}{飽和蒸気圧、湿り蒸気、乾き度、蒸気表、過熱蒸気、モリエ線図、ランキンサイクル、再熱、再生（抽気）、断熱効率、復水器}

\Needspace{11\baselineskip}
\begin{phsaki}{水と蒸気がいっしょに流れる}
この章では、ボイラーや復水器の中を「圧力一定」として、出口と入口の状態だけを数えました。実際のボイラーの管の中では、わいた水と蒸気が混ざって速く流れ、泡の多い所と少ない所ができます。どこで管の壁が乾いて急に熱くなるか（熱-12 の限界熱流束）を予想するのは、原子力発電所の安全にとって大切な問題です。こうした「2つの相がいっしょに流れる」流れ（気液二相流）や、タービンの羽根の形の設計は、蒸気タービン・ボイラー・原子力の熱水力の専門の本で扱います。
\end{phsaki}

%% ---------------- 章末問題 ----------------
\bkendproblems
\bkproblemgroup{基本}{2}
\bkproblem 標高約\qty{2000}{m}の高原の気圧を\qty{80}{kPa}とします。巻末の蒸気表の\qty{90}{\celsius}（\qty{70.182}{kPa}）と\qty{95}{\celsius}（\qty{84.609}{kPa}）の行から、水がわく温度を求めましょう（2けた。\bkreidairef{1}）。
\bkpans{$80-70.182=9.818$、$84.609-70.182=14.427$、$9.818\div14.427\approx0.681$。$90+5\times0.681\approx93.4$ で、約\qty{93}{\celsius}}
\bkproblem \qty{0.2}{MPa}の飽和の状態の、わく温度と蒸発熱を、表\ref{ph:熱-8:tab-sat}から求めましょう（\bkreidairef{2}）。
\bkpans{\qty{120.21}{\celsius}、$L=h''-h'=\qty{2201.6}{kJ/kg}$（巻末の蒸気表の蒸発熱の列と同じ値。表\ref{ph:熱-8:tab-sat}の丸めた値で $2706.2-504.7$ と引くと 2201.5）}
\bkproblem \qty{2}{MPa}・乾き度0.80の湿り蒸気の比エンタルピーを求めましょう（\bkreidairef{4}）。
\bkpans{$2798.4-908.6=1889.8$、$908.6+0.80\times1889.8$ $\approx\qty{2420}{kJ/kg}$}
\bkproblem \qty{1.0}{kW}のヒーターで、1気圧でわいている水を3分間温め続けます。蒸発する水は何 g ですか（\bkreidairef{5}）。
\bkpans{$1.0\times180\div2256.5\approx\qty{0.0798}{kg}$。約\qty{79.8}{g}}
\bkproblem \qty{4}{MPa}・\qty{350}{\celsius}の過熱蒸気の比エンタルピーを、表\ref{ph:熱-8:tab-kanetsu}の\qty{300}{\celsius}と\qty{400}{\celsius}の行から直線補間で求めましょう（\bkreidairef{12}）。
\bkpans{$(2961.7+3214.4)\div2$ $\approx\qty{3088}{kJ/kg}$（巻末の蒸気表では\qty{3093.3}{kJ/kg}。\qty{100}{K}の間かくでは少しずれる）}
\bkproblem ゲージ圧\qty{0.9}{MPa}のボイラー（大気圧\qty{0.1}{MPa}とする）の中の水は、何℃でわいていますか（\bkreidairef{13}）。
\bkpans{絶対圧\qty{1.0}{MPa}。表\ref{ph:熱-8:tab-sat}より約\qty{180}{\celsius}（\qty{179.89}{\celsius}）}
\bkproblem \qty{10}{kPa}・乾き度0.90の湿り蒸気の比体積を、表\ref{ph:熱-8:tab-sat}から求めましょう（\bkreidairef{8}）。
\bkpans{$14.67-0.001010\approx14.669$、$v=0.001010+0.90\times14.669$ $\approx\qty{13.2}{m^3/kg}$}
\bkproblem \bkreidairef{6}のサイクルで、ポンプの仕事は、ボイラーで受けとる熱の何\,\% ですか（\bkreidairef{7}）。
\bkpans{$10.1\div3173.2\approx0.0032$。約0.32\,\%}
\bkproblemgroup{チャレンジ}{4}
\bkproblem \qty{4}{MPa}・\qty{500}{\celsius}の蒸気を\qty{10}{kPa}まで膨張させる理想のランキンサイクルの、タービンの仕事・ポンプの仕事・効率を求めましょう。また、\bkreidairef{6}（\qty{10}{MPa}・\qty{500}{\celsius}）とくらべ、効率と出口の乾き度の違いを説明しましょう（\bkreidairef{6}）。
\bkpans{$x_2=6.4427\div7.4997\approx0.859$、$h_2\approx\qty{2246.8}{kJ/kg}$。$w_{\text{出}}=\qty{1199.0}{kJ/kg}$。$w_{\text{入}}=0.001010\times3990$ $\approx\qty{4.0}{kJ/kg}$。$q_{\text{入}}=3445.8-195.8$ $=\qty{3250.0}{kJ/kg}$。$\eta\approx0.368$。圧力が低いのでわく温度が低く、平均の受熱温度が下がって効率は下がるが、出口の乾き度は上がる}
\bkproblem 冬になって海水が冷たくなり、\bkreidairef{6}の発電所の復水器の圧力が\qty{5}{kPa}に下がりました。\qty{5}{kPa}の飽和表は、$h'=\qty{137.8}{kJ/kg}$、$h''=\qty{2560.8}{kJ/kg}$、$s'=\qty{0.4763}{kJ/(kg.K)}$、$s''=\qty{8.3939}{kJ/(kg.K)}$、$v'=\qty{0.001005}{m^3/kg}$ です。効率と出口の乾き度を求めましょう（\bkreidairef{6}）。
\bkpans{$x_2=6.1230\div7.9176\approx0.773$、$h_2\approx\qty{2011.6}{kJ/kg}$。$w_{\text{出}}=\qty{1363.5}{kJ/kg}$、$w_{\text{入}}\approx\qty{10.0}{kJ/kg}$。$q_{\text{入}}=3375.1-147.8$ $=\qty{3227.3}{kJ/kg}$。$\eta\approx0.419$（\qty{10}{kPa}の0.402より上がる。乾き度は下がる）}
\bkproblem \qty{4}{MPa}・\qty{400}{\celsius}の蒸気を\qty{10}{kPa}まで膨張させるランキンサイクルで、タービンの断熱効率が0.85のとき、タービンの仕事・出口の乾き度・効率を求めましょう（\bkreidairef{6}・じぶんでとこう）。
\bkpans{断熱・可逆なら $x=6.1220\div7.4997$ $\approx0.816$、$h_{2\mathrm s}\approx\qty{2144.5}{kJ/kg}$、差は\qty{1069.9}{kJ/kg}。$w_{\text{出}}=0.85\times1069.9$ $\approx\qty{909.4}{kJ/kg}$。$h_2=\qty{2305.0}{kJ/kg}$、$x_2\approx0.883$。$q_{\text{入}}=\qty{3018.6}{kJ/kg}$、$\eta=905.4\div3018.6$ $\approx0.300$}
\bkproblem \bkreidairef{9}の再熱サイクルで、電気を100万 kW 出すときの蒸気の量と、復水器で捨てる熱を求め、\bkreidairef{11}とくらべましょう（\bkreidairef{9}・\bkreidairef{11}）。
\bkpans{$w_{\text{出}}-w_{\text{入}}=\qty{1461.3}{kJ/kg}$、$\dot m\approx\qty{684}{kg/s}$。$q_{\text{出}}=2246.8-191.8$ $=\qty{2055.0}{kJ/kg}$、捨てる熱 約\qty{1410}{MW}。蒸気は約13\,\%、捨てる熱は約5\,\% 少ない}

\begin{furikaeri}
\begin{checklist}
\item 水のわく温度が、周りの圧力（飽和蒸気圧）で決まることを説明できる。\bkfurref{ph:熱-8:sec1}
\item 乾き度と蒸気表（直線補間）で、湿り蒸気・過熱蒸気の比エンタルピーを求められる。\bkfurref{ph:熱-8:sec2}
\item ランキンサイクルの効率を計算し、再熱・再生の働きを説明できる。\bkfurref{ph:熱-8:sec6}
\end{checklist}
\end{furikaeri}
\bkfukidashi{sora}{水をわかして回す機械の熱の計算ができたね。次の熱-9 では、反対に「わかす」ことで物を冷やす冷蔵庫やエアコンと、空気の中の水蒸気を調べるよ。}
